% ===========================================================================
% appendix_proofs.tex -- the proofs, rewritten line by line (night 5, task G2).
%
% WHAT CHANGED IN G2 AND WHY.  The night-4 skeleton replaced several arguments
% by the phrase "verified symbolically".  A referee cannot check a script, so
% every proof below is now a hand argument that can be read on its own; the
% scripts appear only as an independent second opinion, in a closing line of
% each subsection and in the status comments.  No statement in
% sections/results.tex or sections/model.tex was changed by this task.
%
% STATUS TAGS (CLAUDE.md).  Each subsection carries a comment with the tag of
% the material it contains: [VERIFIED-LP], [VERIFIED-SYMBOLIC],
% [HAND-PROOF-UNREVIEWED], [CONJECTURE].  The visible text never asserts a
% verification status; it states the argument, and the closing line names the
% script that re-checks the identities (or says that no script does).
% The phrases "we prove", "we show" and "it follows" are avoided throughout,
% per CLAUDE.md; so are em dashes, the word "robust", and bare "tight"
% (GLOSSARY.md: "tight" is replaced by "attained with equality" or by a named
% meaning).
%
% ORDERING (conservative choice, G2).  The subsection order of the night-4
% skeleton is preserved verbatim, together with every label
% (app:proofs, app:necessity, app:guarantee, app:tightness, app:coherence,
% app:exact, app:ceiling, app:asymptotics, app:hardness, app:validity,
% app:instances), so that all existing \ref targets keep resolving.  A more
% linear reading order would put app:validity immediately before app:exact,
% since app:exact consumes it; that reordering was NOT made, because it would
% change section numbers that the main text and the night-4 comments refer to.
% Instead app:exact points forward to app:validity explicitly.
%
% Sources reused (with their own status, unchanged):
%   app:tightness .. results/T5_symbolic.py, results/F2_etasel_tight.py
%   app:exact  .... results/N1_dual_certificate.md (dual multipliers),
%                   results/N2_instances.md (attaining instances)
%   app:validity .. results/F2_R6_validity.tex (adapted to the paper notation)
%   app:hardness .. results/N5_bounded_query_hardness.tex
% ===========================================================================

\section{Proofs}\label{app:proofs}

Three conventions are in force throughout this appendix.  First, $f$ is
monotone submodular with $f(\emptyset)=0$ and $\tilde f(\emptyset)=0$, and
submodularity is used in its diminishing-returns form: $d_e(T)\le d_e(S)$ for
$S\subseteq T$ and $e\notin T$.  Second, an optimal set is written $O^{\ast}$,
$|O^{\ast}|\le K$, and instances are normalized so that $f(O^{\ast})=1$
whenever a ratio is computed.  Third, the tie-breaking rule of a greedy run is
part of the instance: "adversarial ties" means that when several elements
maximize the predicted gain, the one named by the construction is taken.

Two of the constructions below are counting functions, so the following
observation is recorded once and used twice.

\begin{lemma}[counting functions]\label{lem:app-count}
Let $N=N_1\cup\dots\cup N_r$ be a partition and let
$h(S)=H(|S\cap N_1|,\dots,|S\cap N_r|)$ for a function $H$ on the box
$\prod_{i}\{0,1,\dots,|N_i|\}$.  Write $\Delta_i H(v)=H(v+u_i)-H(v)$ for the
forward difference in coordinate $i$ ($u_i$ the $i$th unit vector), defined
when $v_i<|N_i|$.  If $\Delta_iH\ge0$ for every $i$ and every $v$, and
$\Delta_\ell\Delta_iH\le0$ for every pair $(i,\ell)$, including $i=\ell$, and
every $v$ at which both differences are defined, then $h$ is monotone and
submodular.
\end{lemma}

\noindent
\emph{Argument.}  Let $v(S)$ denote the count vector of $S$.  For $e\in N_i$
with $e\notin S$ one has $d_e(S)=\Delta_iH(v(S))$, which is nonnegative by
hypothesis, so $h$ is monotone.  For $S\subseteq T$ and $e\in N_i\setminus T$,
the vectors satisfy $v(S)\le v(T)$ coordinatewise and $v(T)_i<|N_i|$.  Walk
from $v(S)$ to $v(T)$ by single unit steps, raising coordinate $i$ last: at
each step the quantity $\Delta_iH$ changes by one second difference
$\Delta_\ell\Delta_iH\le0$, and every intermediate vector has $i$th coordinate
at most $v(T)_i<|N_i|$, so $\Delta_iH$ is defined all along.  Hence
$d_e(T)=\Delta_iH(v(T))\le\Delta_iH(v(S))=d_e(S)$, which is the
diminishing-returns form of submodularity.
% [HAND-PROOF-UNREVIEWED; three lines, no oracle needed.  The lemma is the
% bookkeeping device that turns the (x,y) and (x,z,y) difference tables of
% results/T5_symbolic.py and results/N2_check.py into set-function statements.]

% ---------------------------------------------------------------------------
\subsection{Necessity of an error bound (Proposition~\ref{prop:necessity})}
\label{app:necessity}

Fix $K\ge1$ and $n\ge2K$, and let $\mathcal A$ be an algorithm with arbitrary
query access to $\tilde f$.  Put
\[
  \gamma=\frac{K^{2}}{n(n-K)}\in(0,1],\qquad
  \tilde f(S)=|S| ,
\]
where $\gamma\le1$ because $n\ge2K$ gives $n(n-K)\ge 2K\cdot K>K^{2}$.  For a
$K$-set $O\subseteq N$ define
\[
  f_O(S)=|S\cap O|+\gamma\,|S\setminus O| .
\]

\paragraph{The pair is admissible and has finite error.}
$f_O$ is modular with coefficients in $\{1,\gamma\}\subseteq(0,1]$, so it is
monotone and submodular with $f_O(\emptyset)=0$, and $\tilde f(\emptyset)=0$.
Every single-element gain is $\tilde d_e(S)=1$ and $d_e(S)\in\{1,\gamma\}$, so
$\eta_u=\max_{S,e}d/\tilde d=1$ and $\eta_o=\max_{S,e}\tilde d/d=1/\gamma$: the
error is exactly $\eta=1/\gamma=n(n-K)/K^{2}$, which is finite for every $n$
and grows without bound as $n\to\infty$.  Also
$\max_{|S|\le K}f_O(S)=f_O(O)=K$, because $\gamma\le1$ makes the $K$ largest
modular coefficients exactly those of $O$.

\paragraph{Deterministic algorithms.}
$\tilde f$ does not depend on $O$, so the whole query transcript of
$\mathcal A$, and therefore its output $T$ with $|T|\le K$, are the same for
every $O$.  Since $n\ge2K$ and $|T|\le K$, a $K$-set $O\subseteq N\setminus T$
exists.  For that $O$, $T\cap O=\emptyset$ gives
\[
  \frac{f_O(T)}{f_O(O^{\ast})}\;=\;\frac{\gamma|T|}{K}\;\le\;\gamma
  \;=\;\frac{K^{2}}{n(n-K)}\;\le\;\frac{K}{n-K} .
\]

\paragraph{Randomized algorithms.}
Let $O$ be a uniformly random $K$-subset, independent of the random string of
$\mathcal A$.  The transcript still does not depend on $O$, so the output $T$
is independent of $O$, and the hypergeometric mean gives
$\mathbb E\bigl[|T\cap O|\bigr]=\mathbb E|T|\cdot K/n\le K^{2}/n$.  Bounding
$f_O(T)\le|T\cap O|+\gamma K$ and dividing by $f_O(O^{\ast})=K$,
\[
  \mathbb E\Bigl[\frac{f_O(T)}{f_O(O^{\ast})}\Bigr]
  \;\le\;\frac{K}{n}+\gamma
  \;=\;\frac{K}{n}+\frac{K^{2}}{n(n-K)}
  \;=\;\frac{K(n-K)+K^{2}}{n(n-K)}
  \;=\;\frac{K}{n-K} ,
\]
and $\min_{O}\le\mathbb E_{O}$ produces a single $O$ attaining the bound in
expectation over the random string.  Letting $n\to\infty$ with $K$ fixed drives
the right side to $0$, so no constant worst-case ratio survives when $\eta$ is
unconstrained.
% [HAND-PROOF-UNREVIEWED.]  The night-4 draft used a gamma -> 0 sketch and left
% a note that the constant K/(n-K) of the statement did not match it.  That
% alignment is now settled: the choice gamma = K^2/(n(n-K)) makes the
% RANDOMIZED bound exactly K/(n-K) (the deterministic case is stronger by the
% additive K/n that averaging costs), and n >= 2K is exactly what makes
% gamma <= 1, which is what keeps O optimal.  So the constant in
% Proposition \ref{prop:necessity} is the price of covering randomized
% algorithms, and no statement had to be changed.

No script checks the computations of this proof.

% ---------------------------------------------------------------------------
\subsection{Value accuracy versus gain accuracy
(Proposition~\ref{prop:valueacc})}\label{app:valueacc}
% H1 restoration.  Parts (i) and (iii) are Lemmas 2 and 3 of the TPAMI
% submission (legacy/tpami_submission/appendix.tex), rewritten in the
% notation of Section \ref{sec:model}; part (ii) is a scaling observation
% added in H1.  [HAND-PROOF-UNREVIEWED.]

\paragraph{Part (i).}
Fix $\varepsilon\in(0,1)$ and take $N=\{a,b\}$ with
\[
  f(\emptyset)=0,\quad f(\{a\})=1,\quad
  f(\{b\})=\tfrac{2\varepsilon}{1-\varepsilon},\quad
  f(\{a,b\})=\tfrac{1+\varepsilon}{1-\varepsilon},
\]
\[
  \tilde f(\emptyset)=0,\quad \tilde f(\{a\})=1+\varepsilon,\quad
  \tilde f(\{b\})=\tfrac{2\varepsilon}{1-\varepsilon},\quad
  \tilde f(\{a,b\})=1+\varepsilon .
\]
The function $f$ is nonnegative, monotone, and modular (hence submodular):
$f(\{a,b\})=f(\{a\})+f(\{b\})$ holds because
$1+\tfrac{2\varepsilon}{1-\varepsilon}=\tfrac{1+\varepsilon}{1-\varepsilon}$.
Value accuracy at level $\varepsilon$ is checked set by set: on $\{a\}$,
$\tilde f=(1+\varepsilon)f$; on $\{b\}$, $\tilde f=f$; on $\{a,b\}$,
$\tilde f(\{a,b\})=1+\varepsilon=(1-\varepsilon)\,
\tfrac{1+\varepsilon}{1-\varepsilon}=(1-\varepsilon)f(\{a,b\})$, the lower
end of the band.  At the pair $(S,e)=(\{a\},b)$ the true gain is
$d_b(\{a\})=f(\{a,b\})-f(\{a\})=\tfrac{2\varepsilon}{1-\varepsilon}>0$
while the predicted gain is
$\tilde d_b(\{a\})=\tilde f(\{a,b\})-\tilde f(\{a\})=0$.  A zero predicted
gain at a positive true gain violates
$\tilde d_e(S)\ge d_e(S)/\eta_u$ for every finite $\eta_u$, so no
$(\eta_u,\eta_o)$ of Definition~\ref{def:eta} is finite for this
$\tilde f$, and every bound of the form $L_K(\eta)$ is vacuous for it.

\paragraph{Part (ii).}
Let $f$ be monotone submodular, not identically zero, and
$\tilde f=(1+M)f$ for $M>0$.  Every set $S$ with $f(S)>0$ witnesses
$|\tilde f(S)-f(S)|/f(S)=M>\varepsilon$ for every $\varepsilon<M$, so
$\tilde f$ is not value-accurate at any level below $M$.  Predicted gains
are $\tilde d_e(S)=(1+M)\,d_e(S)$, a positive multiple of the true gains,
so at every state the maximizers of $\tilde d_e(S)$ and of $d_e(S)$ are
the same set of elements, and whichever of them tie breaking picks, the
picked element has maximum true gain.  Every step with positive chosen
gain therefore contributes the factor
$\max_{e\notin S^t}d_e(S^t)/d_{e_t}(S^t)=1$ to
Definition~\ref{def:etasel}, giving $\etasel=1$, and
Proposition~\ref{prop:guarantee} certifies the run at $L_K(1)$.
(Gains are nonnegative by monotonicity, so no step has negative chosen
gain.)

\paragraph{Part (iii).}
Fix $S\subseteq N$ and apply Definition~\ref{def:eta} to the pair
$(\emptyset,S)$ read through single elements: summing the telescoping of
$f$ and $\tilde f$ along any fixed enumeration of $S$, each increment of
$\tilde f$ lies between $1/\eta_u$ and $\eta_o$ times the corresponding
increment of $f$, and increments of a monotone $f$ are nonnegative, so
\[
  f(S)/\eta_u\;\le\;\tilde f(S)\;\le\;\eta_o\,f(S).
\]
Subtracting $f(S)$ and dividing by $f(S)>0$ (the case $f(S)=0$ forces
$\tilde f(S)=0$ and contributes nothing) gives
$\tfrac1{\eta_u}-1\le\tfrac{\tilde f(S)-f(S)}{f(S)}\le\eta_o-1$, hence
$|\tilde f(S)-f(S)|/f(S)\le\max\{1-1/\eta_u,\ \eta_o-1\}$.
% The TPAMI original stated this via d_emptyset(S), i.e. the all-pairs
% form of the error; the single-element form suffices by telescoping, which
% is why the restated part (iii) needs no all-pairs assumption.

No script checks the computations of this proof.

% ---------------------------------------------------------------------------
\subsection{The guarantee (Proposition~\ref{prop:guarantee})}\label{app:guarantee}
This proof is included for completeness; the statement is essentially
Theorem~1 of \citet{goundan2007revisiting}, restated in the prediction-error
model of Section~\ref{sec:model}.
% Decision D2 (TASKS5): attributed, never claimed as new; the words "we prove"
% and "we show" do not appear in this subsection.  [HAND-PROOF-UNREVIEWED.]

Let $S^{0}=\emptyset,S^{1},\dots,S^{K}$ be the states of a run of predictive
greedy with picks $e_{0},\dots,e_{K-1}$ and selection error $\etasel$, and put
$r_t=f(O^{\ast})-f(S^{t})$.

\paragraph{Step 1: one greedy step covers a $1/(\etasel K)$ share of the
residual.}
Let $t<K$ be a step with $d_{e_t}(S^{t})>0$.  Enumerate
$O^{\ast}\setminus S^{t}=\{o_1,\dots,o_m\}$, $m\le|O^{\ast}|\le K$, and
telescope:
\[
  f(O^{\ast}\cup S^{t})-f(S^{t})
  =\sum_{i=1}^{m}\Bigl[f\bigl(S^{t}\cup\{o_1,\dots,o_i\}\bigr)
   -f\bigl(S^{t}\cup\{o_1,\dots,o_{i-1}\}\bigr)\Bigr]
  \;\le\;\sum_{i=1}^{m}d_{o_i}(S^{t}),
\]
one application of submodularity per term.  Monotonicity gives
$f(O^{\ast}\cup S^{t})\ge f(O^{\ast})$, so
$r_t\le\sum_{i=1}^{m}d_{o_i}(S^{t})\le m\max_{e\notin S^{t}}d_e(S^{t})
\le K\max_{e\notin S^{t}}d_e(S^{t})$.  This is the only place where the count
$m\le K$ enters.  At a step with $g_t=d_{e_t}(S^{t})>0$,
Definition~\ref{def:etasel} gives
$M_t=\max_{e\notin S^{t}}d_e(S^{t})=a_tg_t\le\etasel\,d_{e_t}(S^{t})$
whenever $\etasel$ is finite (Step 2 handles the zero steps), hence
\[
  d_{e_t}(S^{t})\;\ge\;\frac{1}{\etasel K}\,r_t,
  \qquad\text{and so}\qquad
  r_{t+1}=r_t-d_{e_t}(S^{t})\;\le\;\Bigl(1-\frac{1}{\etasel K}\Bigr)r_t .
\]

\paragraph{Step 2: the zero-gain steps.}
Monotonicity of $f$ makes $g_t=d_{e_t}(S^{t})\ge0$ always, so the steps not
covered by Step 1 are exactly those with $g_t=0$, and
Definition~\ref{def:etasel} distinguishes two cases.  If the step is harmful
($M_t>0$), then $a_t=\infty$, so $\etasel=\infty$, $L_K(\infty)=0$, and the
claimed bound is the trivial one; nothing needs proving.  If the step is
benign ($M_t=0$), the covering inequality
$r_t\le K\max_{e\notin S^{t}}d_e(S^{t})=KM_t$ of Step 1, whose derivation does
not use positivity of the chosen gain, reads $r_t\le0$: the run has already
reached the optimal value, and $r_s=0$ for every $s\ge t$ by monotonicity.
Therefore, when $\etasel$ is finite, the contraction of Step 1 may be applied
at every step, treating the residual as $0$ from the first zero step onwards.
% [HAND-PROOF-UNREVIEWED.]  Rewritten for the stepwise Definition
% \ref{def:etasel} (decision D3, J2 adoption).  The night-5 version of this
% step argued the harmless case from the GLOBAL band of Definition
% \ref{def:eta}, which is not a hypothesis of the Proposition; the external
% J2 review exhibited a three-element modular counterexample to that reading
% (re-verified, results/H3_j2_recheck.py).  Under the stepwise definition no
% appeal to Definition \ref{def:eta} is needed.  In measured runs on real
% objectives, zero and negative chosen steps are reported separately (columns
% n_steps_nonpos and frac_steps_nonpos of the F1 pipeline).

\paragraph{Step 3: unrolling.}
With $r_0=f(O^{\ast})$, iterating $K$ times gives
$r_K\le(1-\tfrac{1}{\etasel K})^{K}f(O^{\ast})$ and hence
\[
  f(S^{K})=f(O^{\ast})-r_K
  \;\ge\;\Bigl(1-\bigl(1-\tfrac{1}{\etasel K}\bigr)^{K}\Bigr)f(O^{\ast})
  \;=\;L_K(\etasel)\,f(O^{\ast}).
\]
Since $\etasel\ge1$ and $K\ge1$, the quantity $u=1/(\etasel K)$ lies in
$(0,1]$, and $1-u\le e^{-u}$ gives
$L_K(\etasel)\ge1-e^{-1/\etasel}$.

\paragraph{Step 4: the three error measures.}
At every state $S^{t}$ of the run,
\[
  \max_{e\notin S^{t}}d_{e}(S^{t})
  \;\le\;\eta_u\max_{e\notin S^{t}}\tilde d_{e}(S^{t})
  \;=\;\eta_u\,\tilde d_{e_t}(S^{t})
  \;\le\;\eta_u\eta_o\,d_{e_t}(S^{t}),
\]
the first step by the lower band applied to the maximizing element, the middle
equality by the greedy choice, the last by the upper band at $(S^{t},e_t)$.
Reading the two ends as the step ratio $a_t$ of
Definition~\ref{def:etasel} gives, at every step with $g_t>0$,
$a_t\le\eta_u\eta_o$ when the band is the global one, and
$a_t\le\etatr$ when it is restricted to the states of the run; at a step
with $g_t=0$ the same display forces $M_t=0$, so $a_t=1$.  Taking the
maximum over the steps gives $\etasel\le\etatr\le\eta$ for every
predictor with finite error ($\etatr\le\eta$ because the states of the
run are a subset of all sets); this is where the finiteness hypothesis of
the last sentence of the Proposition is used.
The map
$x\mapsto1-\tfrac{1}{xK}$ is increasing on $x\ge1$ with values in $[0,1)$, so
$x\mapsto(1-\tfrac1{xK})^{K}$ is increasing and $L_K$ is decreasing.  Hence
$L_K(\etasel)\ge L_K(\etatr)\ge L_K(\eta)$, and the bound of Step 3 holds a
fortiori with $\etatr$ or $\eta$ in place of $\etasel$.

\begin{remark}[stepwise product bound]\label{rem:app-product}
When the per-step ratios are retained individually, the same unrolling
gives
\[
  f(S^{K})\;\ge\;
  \Bigl(1-\prod_{t=0}^{K-1}\bigl(1-\tfrac1{Ka_t}\bigr)\Bigr)f(O^{\ast}),
  \qquad 1/\infty=0,
\]
which degrades a harmful zero step to a contraction factor of $1$ instead
of collapsing the whole bound to $0$.
% [HAND-PROOF-UNREVIEWED]; suggested by the J2 external review (section
% 4.2), which checked it with exact arithmetic on 1,905 small-instance
% trajectories (results/J1_independent_audit.py); the recursive assembly
% for general runs is not oracle-checked.
\end{remark}

No script certifies the steps of the main argument.  That the bound is met
with equality, and that the three rulers read $\etasel=\etatr=\hat a<\eta$
on the family of Appendix~\ref{app:tightness} (also under the stepwise
Definition~\ref{def:etasel}: no zero step occurs in that family, see
\texttt{results/K1\_etasel\_newdef\_check.py}), is machine-checked by
\texttt{results/F2\_etasel\_tight.py}.

% ---------------------------------------------------------------------------
\subsection{Per-$K$ tightness (Theorem~\ref{thm:tightness})}\label{app:tightness}
% Status: [HAND-PROOF-UNREVIEWED] for the write-up below; every identity in it
% is re-checked symbolically for general K by results/T5_symbolic.py (105/105
% PASS: GEN tier with symbolic K, CONC tier by exhaustive lattice enumeration
% for K = 2..8 in exact arithmetic), and the two run-dependent errors are
% [VERIFIED-LP, K = 2..8, ahat in {1.5, 2}] by results/F2_etasel_tight.py.
% The instance builder used by the experiments is code/check_explicit_instance.py.

Fix $K\ge2$ and $\hat a>1$, and put
\[
  a=1-\frac{1}{\hat aK}\in(0,1),\qquad N=B\cup O,\quad |B|=|O|=K,\quad n=2K .
\]
For $S\subseteq N$ write $x=|S\cap B|$ and $y=|S\cap O|$, and define the pair
$(f,\tilde f)$ by
\begin{equation}\label{eq:uk-family}
\begin{aligned}
  f(S)&=F(x,y)=1-a^{x}\Bigl(1-\frac yK\Bigr),\\[2pt]
  \tilde f(S)&=G(x,y)=
  \begin{cases}
    1-a^{x}, & y=0,\\[2pt]
    1-a^{x}+a^{x}\Bigl[(1-a)+\dfrac{(y-1)a}{K-1}\Bigr], & 1\le y\le K .
  \end{cases}
\end{aligned}
\end{equation}
Both are counting functions of $(x,y)$, so Lemma~\ref{lem:app-count} applies
with $r=2$, $N_1=B$, $N_2=O$.  Note $F(0,0)=0$ and $G(0,0)=0$.

\paragraph{Step 1: a closed form for the second branch of $G$.}
For $1\le y\le K$,
\begin{align*}
  1-a^{x}+a^{x}(1-a)+\frac{(y-1)a^{x+1}}{K-1}
  &=1-a^{x+1}+\frac{(y-1)a^{x+1}}{K-1}\\
  &=1-a^{x+1}\,\frac{(K-1)-(y-1)}{K-1}
   \;=\;1-a^{x+1}\,\frac{K-y}{K-1}.
\end{align*}
In particular $G(x,K)=1$ for every $x$, and the two branches agree at $y=1$:
the second branch gives $1-a^{x+1}$, which is $G(x,0)+a^{x}(1-a)$.

\paragraph{Step 2: the four gain identities.}
Writing $\Delta_x$ and $\Delta_y$ for the forward differences in $x$ and $y$,
the four cases are computed directly from \eqref{eq:uk-family} and Step 1:
\begin{align}
  \Delta_xF(x,y)&=a^{x}(1-a)\,\frac{K-y}{K}, &
  \Delta_yF(x,y)&=\frac{a^{x}}{K},
  \label{eq:uk-F-diffs}\\
  \Delta_xG(x,0)&=a^{x}(1-a), &
  \Delta_yG(x,0)&=a^{x}(1-a),
  \label{eq:uk-G-diffs0}\\
  \Delta_xG(x,y)&=a^{x+1}(1-a)\,\frac{K-y}{K-1}, &
  \Delta_yG(x,y)&=\frac{a^{x+1}}{K-1},
  \label{eq:uk-G-diffs1}
\end{align}
the last row for $1\le y\le K$ in the $x$ direction and for
$1\le y\le K-1$ in the $y$ direction.  The two entries of
\eqref{eq:uk-G-diffs0} are the same number because both read
$(1-a^{x+1})-(1-a^{x})$: the first crosses no branch, and the second crosses
from $G(x,0)$ to $G(x,1)=1-a^{x+1}$ by Step 1.  Dividing
\eqref{eq:uk-G-diffs0} and \eqref{eq:uk-G-diffs1} by the corresponding entries
of \eqref{eq:uk-F-diffs} gives the four ratios, with $R:=aK/(K-1)$:
\[
  \frac{\Delta_xG(x,0)}{\Delta_xF(x,0)}=1,
  \qquad
  \frac{\Delta_yG(x,0)}{\Delta_yF(x,0)}
   =\frac{a^{x}(1-a)}{a^{x}/K}=K(1-a)=\frac1{\hat a},
\]
and, in the range $1\le y\le K-1$ where both are defined,
\[
  \frac{\Delta_xG(x,y)}{\Delta_xF(x,y)}
   =\frac{a^{x+1}(1-a)(K-y)/(K-1)}{a^{x}(1-a)(K-y)/K}=R,
  \qquad
  \frac{\Delta_yG(x,y)}{\Delta_yF(x,y)}
   =\frac{a^{x+1}/(K-1)}{a^{x}/K}=R,
\]
where $K(1-a)=1/\hat a$ is the definition of $a$.  At $y=K$ the only remaining
direction is $x$, and both gains vanish: $\Delta_xF(x,K)=0$ by
\eqref{eq:uk-F-diffs} and $\Delta_xG(x,K)=0$ by Step 1.  So a zero true gain
always comes with a zero predicted gain, as Definition~\ref{def:eta} demands.

\paragraph{Step 3: monotonicity and submodularity of $f$, case by case.}
The two first differences in \eqref{eq:uk-F-diffs} are nonnegative for
$0\le y\le K$ and $x\ge0$, since $a\in(0,1)$.  The four second differences are
\[
  \Delta_x\Delta_xF=-a^{x}(1-a)^{2}\frac{K-y}{K}\le0,\qquad
  \Delta_y\Delta_xF=-\frac{a^{x}(1-a)}{K}\le0,
\]
\[
  \Delta_x\Delta_yF=\frac{a^{x+1}-a^{x}}{K}=-\frac{a^{x}(1-a)}{K}\le0,\qquad
  \Delta_y\Delta_yF=\frac{a^{x}}{K}-\frac{a^{x}}{K}=0\le0 .
\]
The two mixed differences coincide, as they must.  Lemma~\ref{lem:app-count}
now gives that $f$ is monotone and submodular on the whole of $2^{N}$.  (No
claim is made about $\tilde f$; it is monotone but not submodular, which the
model of Section~\ref{sec:model} permits.)

\paragraph{Step 4: the global error of the pair.}
By Step 2 the set of single-element ratios $\tilde d_e(S)/d_e(S)$, over all
$S$ and all $e\notin S$ with $d_e(S)>0$, is exactly $\{1,\ R,\ 1/\hat a\}$.
Two sign computations order it:
\[
  R-1=\frac{aK-(K-1)}{K-1}=\frac{K-\tfrac1{\hat a}-K+1}{K-1}
     =\frac{1-1/\hat a}{K-1}>0,\qquad
  1-\frac1{\hat a}=\frac{\hat a-1}{\hat a}>0 ,
\]
both strict because $\hat a>1$.  Hence
$\eta_o=\max\{1,R,1/\hat a\}=R$ and $\eta_u=\max\{1,1/R,\hat a\}=\hat a$, both
attained, and
\begin{equation}\label{eq:uk-eta}
  \eta=\eta_u\eta_o=\hat a\,\frac{aK}{K-1}
   =\frac{K(\hat a-\tfrac1K)}{K-1}=\frac{K\hat a-1}{K-1},
\end{equation}
using $\hat aa=\hat a-1/K$.

\paragraph{Step 5: the optimum.}
$1-F(x,y)=a^{x}(K-y)/K\ge0$, so $F\le1$ everywhere, and $F(0,K)=1$.  Since
$|O|=K$, the value $f(O^{\ast})=1$ is attained at $O$, and $O^{\ast}=O$ may be
taken.

\paragraph{Step 6: the greedy induction under adversarial ties.}
Claim: for $0\le t\le K$ the state after $t$ steps satisfies $S^{t}\subseteq B$
and $|S^{t}|=t$, that is $(x,y)=(t,0)$.  For $t=0$ this is $S^{0}=\emptyset$.
Assume it at $t<K$.  At the state $(t,0)$ every candidate is either a
remaining element of $B$, whose predicted gain is $\Delta_xG(t,0)=a^{t}(1-a)$,
or an element of $O$, whose predicted gain is $\Delta_yG(t,0)=a^{t}(1-a)$;
these are equal by \eqref{eq:uk-G-diffs0}, so every candidate is a maximizer
and the step is a tie among all $2K-t$ of them.  Adversarial tie breaking
selects an element of $B$ (which is possible because $t<K=|B|$), so
$S^{t+1}$ has $(x,y)=(t+1,0)$.  After $K$ steps $S^{K}=B$ and
\[
  f(S^{K})=F(K,0)=1-a^{K}
  =1-\Bigl(1-\frac{1}{\hat aK}\Bigr)^{K}=L_K(\hat a),
\]
which is $L_K(\hat a)f(O^{\ast})$ by Step 5.
% [HAND-PROOF-UNREVIEWED: the induction itself.  Its per-step content is the
% tie identity (item 4 of results/T5_symbolic.py, verified for symbolic K), and
% the full simulation with symbolic tie comparison is re-run for K = 2..8 in
% the CONC tier of the same script.]

\paragraph{Step 7: the two run-dependent errors equal $\hat a$.}
At the state $S^{t}=(t,0)$ the true gains are $a^{t}(1-a)=a^{t}/(\hat aK)$ for
an element of $B$ and $a^{t}/K$ for an element of $O$, by
\eqref{eq:uk-F-diffs}.  Since $\hat a>1$, the largest true gain at that state
is $a^{t}/K$ while the chosen element has gain $a^{t}/(\hat aK)>0$, so every
step has positive chosen gain, the stepwise $a_t$ of
Definition~\ref{def:etasel} reduce to the ratio form, and
\[
  a_t=\frac{\max_{e\notin S^{t}}d_e(S^{t})}{d_{e_t}(S^{t})}
  =\frac{a^{t}/K}{a^{t}/(\hat aK)}=\hat a
\]
at every one of the $K$ steps, so $\etasel=\hat a$.
% D3 check (J2 adoption): no zero step occurs in this family, so the
% stepwise definition changes nothing here; machine-checked for K=2..8,
% ahat in {1.5,2} by results/K1_etasel_newdef_check.py (14/14).  For $\etatr$, only the
states of the run matter, and at those states $y=0$, so by Step 2 the only two
ratios that occur are $1$ (direction $x$) and $1/\hat a$ (direction $y$);
therefore the trajectory-restricted factors are
$\eta_o^{\mathrm{tr}}=\max\{1,1/\hat a\}=1$ and
$\eta_u^{\mathrm{tr}}=\max\{1,\hat a\}=\hat a$, so $\etatr=\hat a$ as well.
Comparing with \eqref{eq:uk-eta},
$\eta-\hat a=\frac{K\hat a-1-(K-1)\hat a}{K-1}=\frac{\hat a-1}{K-1}>0$, so on
this family the selection and trajectory rulers agree and are strictly smaller
than the global one.

\paragraph{Step 8: the same family read at fixed global error.}
Solving \eqref{eq:uk-eta} for $\hat a$ gives $\hat a=(\eta(K-1)+1)/K$, hence
$a=1-\frac{1}{\hat aK}=1-\frac{1}{\eta(K-1)+1}$ and
\[
  f(S^{K})=1-a^{K}
  =1-\Bigl(1-\frac{1}{\eta(K-1)+1}\Bigr)^{K}=U_K(\eta).
\]
So the same instances that make Proposition~\ref{prop:guarantee} an equality at
$\etasel=\hat a$ realize only $U_K$ when charged the global error $\eta$, which
is the comparison recorded in Remark~\ref{rem:exact-gap}.
% Together with Appendix \ref{app:exact} this is where U_K and min_j V_j part
% company: U_K - V_{K-1} = q^{K-1}(eta-1)/(K eta k_1) > 0 for eta > 1
% [VERIFIED-SYMBOLIC, results/N1_dual_certificate.py part D].

Steps 1 to 8 give Theorem~\ref{thm:tightness}: for every $K\ge2$ and every
$\hat a>1$ the pair \eqref{eq:uk-family} on $2K$ elements has monotone
submodular $f$, finite error, an adversarial-tie greedy run with
$\etasel=\etatr=\hat a$, and output value exactly $L_K(\hat a)f(O^{\ast})$.

The identities of this proof are machine-checked by
\texttt{results/T5\_symbolic.py} and \texttt{results/F2\_etasel\_tight.py}.

% ---------------------------------------------------------------------------
\subsection{Coherence (Lemma~\ref{lem:coherence})}\label{app:coherence}
% Status: [HAND-PROOF-UNREVIEWED]; four lines, single-element bounds only.
% Renamed from "consistency lemma": consistency keeps its
% learning-augmented-algorithms meaning throughout (GLOSSARY.md).

Let $S\subseteq N$ and $e,e'\notin S$ with $e\ne e'$, and assume
$\tilde d_{e}(S)\ge\tilde d_{e'}(S)$.

\paragraph{The exchange identity.}
Expanding the same quantity $\tilde f(S\cup\{e,e'\})-\tilde f(S)$ in the two
possible orders, first $e$ then $e'$, and first $e'$ then $e$:
\[
  \tilde d_{e}(S)+\tilde d_{e'}(S\cup\{e\})
  \;=\;\tilde f(S\cup\{e,e'\})-\tilde f(S)
  \;=\;\tilde d_{e'}(S)+\tilde d_{e}(S\cup\{e'\}) .
\]
This uses nothing but the definition of a set function.  Subtracting and
applying the hypothesis $\tilde d_{e}(S)-\tilde d_{e'}(S)\ge0$,
\begin{equation}\label{eq:coh-pred}
  \tilde d_{e'}(S\cup\{e\})-\tilde d_{e}(S\cup\{e'\})
  =\tilde d_{e'}(S)-\tilde d_{e}(S)\;\le\;0 .
\end{equation}

\paragraph{Part (i).}
Chain the two bands of Definition~\ref{def:eta} around
\eqref{eq:coh-pred}:
\[
  d_{e'}(S\cup\{e\})\;\le\;\eta_u\,\tilde d_{e'}(S\cup\{e\})
  \;\le\;\eta_u\,\tilde d_{e}(S\cup\{e'\})
  \;\le\;\eta_u\eta_o\,d_{e}(S\cup\{e'\})
  \;=\;\eta\,d_{e}(S\cup\{e'\}),
\]
the first inequality by the lower band at $(S\cup\{e\},e')$, the second by
\eqref{eq:coh-pred}, the third by the upper band at $(S\cup\{e'\},e)$.
Dividing by $\eta$ gives (i).

\paragraph{Part (ii).}
The same two-order expansion applied to $f$ reads
$d_{e}(S)+d_{e'}(S\cup\{e\})=d_{e'}(S)+d_{e}(S\cup\{e'\})$, that is
\[
  d_{e'}(S)-d_{e}(S)\;=\;d_{e'}(S\cup\{e\})-d_{e}(S\cup\{e'\}) .
\]
Substituting part (i) in the form
$-d_{e}(S\cup\{e'\})\le-\tfrac1\eta d_{e'}(S\cup\{e\})$,
\[
  d_{e'}(S)-d_{e}(S)\;\le\;d_{e'}(S\cup\{e\})
   -\frac1\eta\,d_{e'}(S\cup\{e\})
  \;=\;\Bigl(1-\frac1\eta\Bigr)d_{e'}(S\cup\{e\}),
\]
which is (ii).  Only $\eta\ge1$ is used beyond the bands, and only to keep the
factor $1-\tfrac1\eta$ nonnegative when the inequality is applied later.

No script checks the computations of this proof.  The two inequalities are
used as constraint families in Appendix~\ref{app:validity}, and the linear
program they define is compared against the full-lattice program by
\texttt{results/N3\_K5\_lattice.py}.

% ---------------------------------------------------------------------------
\subsection{The exact worst case (Theorem~\ref{thm:exact})}\label{app:exact}

\paragraph{Roadmap.}
The mechanism behind the value runs in the following order.  The sharp form of
Lemma~\ref{lem:coherence}, in Section~\ref{sec:coherence}, limits how much a
single step can lose now and cost later at the same time; the two-phase
trajectory of the attaining instances below is the shape that limit leaves
available; the multipliers \eqref{eq:duals-jpos} turn that shape into a lower
bound valid for every run, and the instances \eqref{eq:vj-family} attain it;
Appendix~\ref{app:rigidity} reads the same certificate backwards and recovers
the trajectory from the equality case; and the integer breakpoints appear
there as the error levels at which the multipliers carrying that argument
vanish.  The coefficient bookkeeping of the certificate, the gain table of the
instances, and the validity of the four constraint families
(Appendix~\ref{app:validity}) are kept apart from that line of reasoning and
can be read afterwards.
% Roadmap added in H-J3 item 5, following results/J3_proof_structure.md
% section 5.  Minimal reorganization: the existing proofs below are unchanged,
% only this pointer paragraph was inserted.

Throughout this subsection $K\ge2$, $\eta\ge1$, $f(O^{\ast})=1$, and
\[
  k_1=(K-1)\eta+1,\qquad q=\frac{(K-1)\eta}{k_1}=1-\frac1{k_1},\qquad
  V_j(\eta)=1-q^{\,j}\Bigl(1-\frac{K-j}{K\eta}\Bigr).
\]
Note $k_1\ge K\ge2$, so $q\in(0,1)$, and $(1-q)k_1=1$.

\paragraph{The reduced linear program.}
Every execution of predictive greedy induces $d_t=d_{e_t}(S^{t})$ for
$t=0,\dots,K-1$ and, after enumerating $O^{\ast}=\{o_1,\dots,o_K\}$,
\[
  g_{t,i}=\begin{cases} d_{o_i}(S^{t}), & o_i\notin S^{t},\\ 0,& o_i\in S^{t},
  \end{cases}\qquad 0\le t\le K,\ 1\le i\le K .
\]
Appendix~\ref{app:validity} derives, one family at a time and with the case
$e_t\in O^{\ast}$ treated explicitly, that this vector satisfies
\begin{equation}\label{eq:redlp}
\begin{aligned}
  \textstyle\sum_{i}g_{t,i}+\sum_{s<t}d_s&\ \ge\ 1 && (\mathrm{sum}(t)),\\
  d_t-g_{t,i}/\eta&\ \ge\ 0 && (\mathrm{pred}(t,i)),\\
  g_{t,i}-g_{t+1,i}&\ \ge\ 0 && (\mathrm{mono}(t,i)),\\
  \bigl(1-\tfrac1\eta\bigr)g_{t+1,i}-g_{t,i}+d_t&\ \ge\ 0
      && (\mathrm{cons}(t,i)),
\end{aligned}
\end{equation}
together with $d_t\ge0$, $g_{t,i}\ge0$, and that $\sum_{t<K}d_t=f(S^{K})$ by
telescoping.  Consequently the minimum of $\sum_{t<K}d_t$ over the polyhedron
\eqref{eq:redlp} is at most the worst-case ratio $\rho_K(\eta)$: the
polyhedron contains every realizable run and possibly more.

\paragraph{Lower bound: greedy is never worse than $\min_jV_j$.}
Fix a segment index $j$ with $0\le j\le K-1$ and let $\eta$ lie in
$[K-j,K-j+1]$ (in $[K,\infty)$ when $j=0$).  Write $M=K\eta-(K-j)$ and define
nonnegative numbers $\lambda_S(t)$, $\lambda_P(t)$, $\lambda_C(t)$ (one per
constraint family, the same for every index $i$ by the symmetry of
\eqref{eq:redlp} in $i$) as follows.  For $j\ge1$,
\begin{equation}\label{eq:duals-jpos}
\begin{aligned}
  \lambda_S(0)&=\frac{q^{\,j-1}M}{Kk_1}, &
  \lambda_S(t)&=\frac{q^{\,j-1-t}M}{k_1^{2}}\quad(1\le t\le j-1),\\
  \lambda_S(j)&=\frac{K-j+1-\eta}{k_1}, &
  \lambda_C(t)&=\frac{q^{\,j-1-t}M}{Kk_1}\quad(t\le j-1),\\
  \lambda_C(t)&=\frac{K-1-t}{K(\eta-1)}\quad(t\ge j), &
  \lambda_P(t)&=\frac{\eta-(K-t)}{K(\eta-1)}\quad(t\ge j),
\end{aligned}
\end{equation}
with $\lambda_S(t)=0$ for $t>j$ and $\lambda_P(t)=0$ for $t<j$; for $j=0$,
\begin{equation}\label{eq:duals-jzero}
\begin{aligned}
  \lambda_S(0)&=\frac1\eta,\qquad \lambda_S(t)=0\quad(t\ge1),\\
  \lambda_C(t)&=\frac{K-1-t}{K(\eta-1)},\qquad
  \lambda_P(t)=\frac{\eta-(K-t)}{K(\eta-1)}\qquad(0\le t\le K-1).
\end{aligned}
\end{equation}
All multipliers of the family $\mathrm{mono}$ are $0$: the lower bound does not
use that the optimal gains decrease along the run.
% [VERIFIED-SYMBOLIC, general symbolic (K, j, t, i) modulo one finite branch
% enumeration, plus independent brute force for K = 2..10:
% results/N1_dual_certificate.py, 320/320 PASS.  Corollary A there: deleting
% the mono rows does not change the LP value.]

\emph{Nonnegativity.}  On the segment, $\eta\ge K-j\ge1$ and $K\eta>K-j$, so
$M>0$; since also $q\in(0,1)$ and $k_1>0$, the three multipliers built from
powers of $q$, namely $\lambda_S(0)$, $\lambda_S(t)$ for $1\le t\le j-1$ and
$\lambda_C(t)$ for $t\le j-1$, are positive.  For
$\lambda_S(j)$, the segment condition $\eta\le K-j+1$ is exactly
$K-j+1-\eta\ge0$, and equality holds at the right endpoint of the segment.
For $\lambda_C(t)$ with $t\ge j$, $K-1-t\ge0$ because $t\le K-1$.  For
$\lambda_P(t)$ with $t\ge j$, $\eta-(K-t)\ge\eta-(K-j)\ge0$ by the segment
condition, with equality at $t=j$ at the left endpoint.  The denominator
$\eta-1$ is positive whenever $j\le K-2$, since then $\eta\ge K-j\ge2$.  The
one remaining case is $j=K-1$, whose segment is $[1,2]$: there the range
$t\ge j$ contains only $t=K-1$, where the numerator $K-1-t$ vanishes
identically and $\lambda_C(K-1)$ is defined to be $0$, while
$\lambda_P(K-1)=(\eta-1)/(K(\eta-1))$ is defined by cancellation to be $1/K$.
So \eqref{eq:duals-jpos} has no singularity on any segment.

\emph{The weighted-sum identity.}  For an arbitrary vector $(d_t,g_{t,i})$,
not necessarily feasible, consider
\begin{equation}\label{eq:combination}
\begin{aligned}
  \Sigma\;:=\;&\sum_{t=0}^{j}\lambda_S(t)
      \Bigl[\sum_ig_{t,i}+\sum_{s<t}d_s-1\Bigr]
   \;+\;\sum_{t=j}^{K-1}\sum_i\lambda_P(t)
      \Bigl[d_t-\frac{g_{t,i}}{\eta}\Bigr]\\
  &+\sum_{t=0}^{K-1}\sum_i\lambda_C(t)
     \Bigl[\bigl(1-\tfrac1\eta\bigr)g_{t+1,i}-g_{t,i}+d_t\Bigr].
\end{aligned}
\end{equation}
The last summand of the third group is void, because $\lambda_C(K-1)=0$ by its
own formula in \eqref{eq:duals-jpos} and \eqref{eq:duals-jzero}.
Collecting coefficients in $\Sigma$ gives three families of scalar identities.
The coefficient of $g_{t,i}$, for $0\le t\le K$, is
$\lambda_S(t)-\lambda_P(t)/\eta-\lambda_C(t)+(1-\tfrac1\eta)\lambda_C(t-1)$,
with the conventions $\lambda_C(-1)=0$, $\lambda_C(K)=0$ and
$\lambda_S(t)=\lambda_P(t)=0$ outside the ranges listed above; the coefficient
of $d_t$ is $\sum_{s>t}\lambda_S(s)+K\bigl(\lambda_P(t)+\lambda_C(t)\bigr)$;
the constant is $-\sum_{t\le j}\lambda_S(t)$.  The claim is that these equal
$0$, $1$ and $-V_j$ respectively, so that
$\Sigma=\sum_{t<K}d_t-V_j$ identically.  Each is checked below; nothing is
left to "direct expansion".

\emph{(a) The coefficient of $g_{t,i}$ vanishes.}  Take $j\ge1$ first.
For $t=0$ one has $\lambda_C(-1)=0$ and $\lambda_P(0)=0$ (as $0<j$), so the
coefficient is $\lambda_S(0)-\lambda_C(0)=0$, the two entries of
\eqref{eq:duals-jpos} being literally the same number.  For $1\le t\le j-1$,
again $\lambda_P(t)=0$, and dividing the coefficient by the positive quantity
$q^{\,j-1-t}M/(Kk_1)$ turns it into
\[
  \frac{K}{k_1}-1+\Bigl(1-\frac1\eta\Bigr)q
  =\frac{K-k_1}{k_1}+\frac{\eta-1}{\eta}\cdot\frac{(K-1)\eta}{k_1}
  =\frac{(K-1)(1-\eta)+(\eta-1)(K-1)}{k_1}=0,
\]
using $K-k_1=K-(K-1)\eta-1=(K-1)(1-\eta)$.  For $t=j$, write $u=K-j$, so that
$\eta\in[u,u+1]$ and $M=K\eta-u$, and multiply the coefficient by
$K\eta k_1(\eta-1)>0$:
\[
  \underbrace{K\eta(\eta-1)(u+1-\eta)}_{\text{from }\lambda_S(j)}
  \underbrace{-\,k_1(\eta-u)}_{\text{from }\lambda_P(j)/\eta}
  \underbrace{-\,(u-1)\eta k_1}_{\text{from }\lambda_C(j)}
  \underbrace{+\,(\eta-1)^{2}(K\eta-u)}_{\text{from }(1-1/\eta)\lambda_C(j-1)} .
\]
The two middle terms combine as
$-k_1\bigl[(\eta-u)+(u-1)\eta\bigr]=-k_1u(\eta-1)$, so with $A=\eta-1$ the
whole expression is $A\bigl[K\eta(u+1-\eta)+A(K\eta-u)-k_1u\bigr]$, and the
bracket expands to
\[
  K\eta u+K\eta-K\eta^{2}+K\eta^{2}-K\eta-\eta u+u-k_1u
  =u\bigl[K\eta-\eta+1-k_1\bigr]=u\,[\,k_1-k_1\,]=0 .
\]
For $j+1\le t\le K-1$ one has $\lambda_S(t)=0$, and multiplying by
$K\eta(\eta-1)$ the coefficient becomes
\[
  -(\eta-K+t)-\eta(K-1-t)+(\eta-1)(K-t)
  =-\eta+K-t-\eta K+\eta+\eta t+\eta K-\eta t-K+t=0 .
\]
For $t=K$ only the last term survives and $\lambda_C(K-1)=0$.  The
multiplication by $\eta-1$ in the last two computations is legitimate on every
segment with $j\le K-2$, where $\eta\ge2$; on the segment $j=K-1$, which is
$[1,2]$, both sides of each identity are rational functions of $\eta$ that are
continuous at $\eta=1$ after the cancellation recorded above, so the identity
at $\eta=1$ is the limit of the identity on $(1,2]$.  For $j=0$ the
first computation is replaced by the single case $t=0$: multiplying by
$K\eta(\eta-1)$,
$K(\eta-1)-(\eta-K)-\eta(K-1)=K\eta-K-\eta+K-\eta K+\eta=0$, and the cases
$1\le t\le K-1$ and $t=K$ are the ones just done.

\emph{(b) The coefficient of $d_t$ equals $1$.}  For $t\ge j$ the sum
$\sum_{s>t}\lambda_S(s)$ is empty, and
\[
  K\bigl(\lambda_P(t)+\lambda_C(t)\bigr)
  =K\cdot\frac{(\eta-K+t)+(K-1-t)}{K(\eta-1)}=\frac{\eta-1}{\eta-1}=1 ,
\]
which at $j=K-1$, $\eta=1$ reads $K(\tfrac1K+0)=1$ with the cancelled values.
For $t=j-1$ (so $j\ge1$), $\lambda_P(j-1)=0$ and the only surviving
$\lambda_S$ is $\lambda_S(j)$, so the coefficient is
\[
  \frac{K-j+1-\eta}{k_1}+K\cdot\frac{M}{Kk_1}
  =\frac{(K-j+1-\eta)+(K\eta-K+j)}{k_1}=\frac{(K-1)\eta+1}{k_1}=1 .
\]
For $t\le j-2$ the value is unchanged from $t+1$ to $t$: subtracting the two
coefficients and using $\lambda_P(t)=\lambda_P(t+1)=0$,
\begin{align*}
  \lambda_S(t+1)+K\bigl[\lambda_C(t)-\lambda_C(t+1)\bigr]
  &=\frac{q^{\,j-2-t}M}{k_1^{2}}
   +K\cdot\frac{M q^{\,j-2-t}}{Kk_1}\,(q-1)\\
  &=\frac{q^{\,j-2-t}M}{k_1^{2}}-\frac{Mq^{\,j-2-t}}{k_1^{2}}=0,
\end{align*}
because $q-1=-1/k_1$.  Downward induction from $t=j-1$ therefore gives the
value $1$ for every $t\le j-1$ as well.

\emph{(c) The constant equals $-V_j$.}  For $j\ge1$, the middle multipliers
form a geometric sum with ratio $q$, which is summed with $1-q=1/k_1$:
\[
  \sum_{t=1}^{j-1}\lambda_S(t)=\frac{M}{k_1^{2}}\sum_{m=0}^{j-2}q^{m}
  =\frac{M}{k_1^{2}}\cdot\frac{1-q^{\,j-1}}{1-q}
  =\frac{M\bigl(1-q^{\,j-1}\bigr)}{k_1},
\]
so that, using $M+(K-j)=K\eta$ and hence $M+(K-j+1-\eta)=K\eta-\eta+1=k_1$,
\[
  \sum_{t=0}^{j}\lambda_S(t)
  =\frac{1}{k_1}\Bigl[\frac{Mq^{\,j-1}}{K}+M-Mq^{\,j-1}+(K-j+1-\eta)\Bigr]
  =\frac{1}{k_1}\Bigl[k_1-Mq^{\,j-1}\frac{K-1}{K}\Bigr].
\]
The definition of $q$ gives $\frac{K-1}{K}=\frac{qk_1}{K\eta}$, so
$Mq^{\,j-1}\frac{K-1}{K}=\frac{Mq^{\,j}k_1}{K\eta}$ and
\[
  \sum_{t=0}^{j}\lambda_S(t)=1-\frac{Mq^{\,j}}{K\eta}
  =1-q^{\,j}\,\frac{K\eta-(K-j)}{K\eta}
  =1-q^{\,j}\Bigl(1-\frac{K-j}{K\eta}\Bigr)=V_j .
\]
For $j=0$ the sum is the single term $\lambda_S(0)=1/\eta$, and
$V_0=1-(1-\tfrac{K}{K\eta})=1/\eta$.

\emph{Conclusion of the lower bound.}  By (a), (b), (c), identity
\eqref{eq:combination} reads $\sum_{t<K}d_t-V_j=\Sigma$ for every vector
$(d_t,g_{t,i})$.  At a point feasible for \eqref{eq:redlp} every bracket in
$\Sigma$ is nonnegative and every multiplier is nonnegative, so $\Sigma\ge0$
and $\sum_{t<K}d_t\ge V_j$.  Since every run of predictive greedy at error
$\eta$ gives such a feasible point with objective $f(S^{K})/f(O^{\ast})$,
\[
  \rho_K(\eta)\;\ge\;V_j(\eta)\qquad\text{on the segment of $j$.}
\]

\emph{Which $V_j$ is smallest.}  For $0\le i\le K-2$,
% index upper limit corrected from K-1 in K6 (J2 review section 3): i = K-1
% would invoke the undefined V_K; the comparisons below only ever use
% adjacent pairs up to (V_{K-2}, V_{K-1}).
\begin{align*}
  V_i-V_{i+1}
  &=q^{i}\Bigl[q-\frac{q(K-i-1)}{K\eta}-1+\frac{K-i}{K\eta}\Bigr]\\
  &=q^{i}\Bigl[-\frac1{k_1}+\frac{1}{K\eta}
    \Bigl(1+\frac{K-i-1}{k_1}\Bigr)\Bigr]
  \;=\;\frac{q^{i}\,(K-i-\eta)}{K\eta k_1},
\end{align*}
where the middle step uses $q-1=-1/k_1$ and
$(K-i)-q(K-i-1)=1+\frac{K-i-1}{k_1}$, and the last step uses
$k_1-K\eta=1-\eta$.  On the segment $\eta\in[K-j,K-j+1]$ the sign of
$K-i-\eta$ is nonnegative for $i\le j-1$ and nonpositive for $i\ge j$, so the
sequence $V_0,\dots,V_{K-1}$ decreases up to index $j$ and increases after it:
$\min_iV_i=V_j$ there, and neighbouring segments agree at the integer
endpoints.  The breakpoints are therefore the integers $2,\dots,K$, and
$\rho_K(\eta)\ge\min_iV_i(\eta)$ for every $\eta\ge1$.  At $\eta\ge K$ the
active index is $j=0$ and $\min_iV_i=V_0=1/\eta$.
The closed forms displayed in Theorem~\ref{thm:exact} for $K\in\{2,3,4\}$ are
the specializations of $\min_jV_j$; for instance at $K=3$ and $j=2$,
$V_2=1-\frac{4\eta^{2}}{(2\eta+1)^{2}}\cdot\frac{3\eta-1}{3\eta}
=\frac{16\eta+3}{3(2\eta+1)^{2}}$.

The identities of this proof are machine-checked by
\texttt{results/N1\_dual\_certificate.py}, and the $K\in\{2,3,4\}$ closed forms
by \texttt{results/T3\_K3\_closed\_form.py}.

\paragraph{Upper bound: the value $V_j$ is attained.}
Fix $K\ge2$, $0\le j\le K-1$, and $\eta$ with $\eta\ge K-j$, and fix any split
$\eta_u\eta_o=\eta$ with $\eta_u,\eta_o\ge1$ (the scripts use
$\eta_u=\eta_o=\sqrt\eta$).  Partition a ground set of $n=2K$ elements into
\[
  C\ (|C|=j),\qquad P\ (|P|=K-j),\qquad O\ (|O|=K),
\]
write $x=|S\cap C|$, $z=|S\cap P|$, $y=|S\cap O|$, and put
\[
  \delta_j=\frac{q^{\,j}}{K\eta},\qquad
  \tilde\delta_j=\eta_o\delta_j=\frac{q^{\,j}}{K\eta_u},\qquad
  W(0)=\frac{k_1}{K\eta_u},\qquad
  W(y)=\frac{(K-y)\eta_o}{K}\ \ (1\le y\le K),
\]
and $\chi(y)=\mathbf 1[y<K]$.  The instance is
\begin{equation}\label{eq:vj-family}
  f(S)=1-q^{x}\Bigl(1-\frac yK\Bigr)+z\,\delta_j\,\chi(y),
  \qquad
  \tilde f(S)=W(0)-q^{x}W(y)+z\,\eta_o\delta_j\,\chi(y).
\end{equation}
Both are counting functions of $(x,z,y)$, so Lemma~\ref{lem:app-count} applies
with $r=3$.  Two normalizations are immediate: $f(\emptyset)=1-1=0$ and
$\tilde f(\emptyset)=W(0)-W(0)=0$.
% Instance family of results/N2_instances.md section 1 (capped variant, the
% script default).  The uncapped variant chi == 1 is monotone submodular for
% every eta >= 1 but is not needed here: the greedy run never leaves y = 0.

\emph{Marginal gains.}  Directly from \eqref{eq:vj-family}, for $y<K$,
\[
  \Delta_Cf=q^{x}(1-q)\frac{K-y}{K},\qquad
  \Delta_Pf=\delta_j,\qquad
  \Delta_Of=\begin{cases}q^{x}/K,& y\le K-2,\\
                        q^{x}/K-z\delta_j,& y=K-1,\end{cases}
\]
and at $y=K$ all three gains are $0$.  On the predictor side, $\Delta_C\tilde f=q^{x}(1-q)W(y)$ and
$\Delta_P\tilde f=\eta_o\delta_j$, while
\[
  \Delta_O\tilde f=\begin{cases}
    q^{x}\bigl(W(0)-\tfrac{(K-1)\eta_o}{K}\bigr), & y=0,\\
    q^{x}\eta_o/K, & 1\le y\le K-2,\\
    q^{x}\eta_o/K-z\eta_o\delta_j, & y=K-1 .
  \end{cases}
\]
Two identities are used repeatedly: $(1-q)k_1=1$ and
\begin{equation}\label{eq:W0-identity}
  W(0)-\frac{(K-1)\eta_o}{K}
  =\frac{(K-1)\eta+1}{K\eta_u}-\frac{(K-1)\eta_u\eta_o}{K\eta_u}
  =\frac{1}{K\eta_u}.
\end{equation}

\emph{Monotonicity and submodularity of $f$.}  The three first differences are
nonnegative: $\Delta_Cf\ge0$ and $\Delta_Pf>0$ are immediate from
$q\in(0,1)$, and at $y=K-1$
\[
  \Delta_Of=\frac{q^{x}}{K}-z\,\frac{q^{\,j}}{K\eta}\ \ge\ 0
  \iff \eta\,q^{\,x-j}\ \ge\ z ,
\]
which holds because $x\le j$ makes $q^{\,x-j}\ge1$ and $z\le K-j\le\eta$ by the
standing assumption $\eta\ge K-j$.  This is the only place where the segment
condition is consumed, and it is where the breakpoints come from.  The nine
second differences are:
$\Delta_C\Delta_Cf=-q^{x}(1-q)^{2}\frac{K-y}{K}\le0$;
$\Delta_O\Delta_Cf=-q^{x}(1-q)/K\le0$ for every $y\le K-1$ (at $y=K-1$ the
gain drops from $q^{x}(1-q)/K$ to $0$, the same decrement);
$\Delta_C\Delta_Of=-q^{x}(1-q)/K\le0$;
$\Delta_P\Delta_Cf=\Delta_C\Delta_Pf=\Delta_P\Delta_Pf=0$;
$\Delta_O\Delta_Pf=0$ for $y\le K-2$ and $-\delta_j\le0$ at $y=K-1$;
$\Delta_P\Delta_Of=0$ for $y\le K-2$ and $-\delta_j\le0$ at $y=K-1$;
$\Delta_O\Delta_Of=0$ for $y\le K-3$ and $-z\delta_j\le0$ at $y=K-2$.
Lemma~\ref{lem:app-count} then gives that $f$ is monotone and submodular.
% The full list of nine second differences with symbolic K is checked in Part C
% of results/N2_check.py (44 symbolic items), and the whole 2^{2K} lattice is
% enumerated for K = 2..6 in Part A.  ftilde is monotone but not submodular
% (Part G): the model of Section \ref{sec:model} does not require it to be.

\emph{The error is exactly $(\eta_u,\eta_o)$.}  Divide each predicted gain by
the corresponding true gain, using \eqref{eq:W0-identity} in the case $y=0$ of
direction $O$.  Every single-element ratio $\tilde d/d$ takes one of three
values:
\[
  \frac{\Delta_C\tilde f}{\Delta_Cf}\Big|_{y=0}=W(0)=\frac{k_1}{K\eta_u},
  \qquad
  \frac{\Delta_O\tilde f}{\Delta_Of}\Big|_{y=0}=\frac1{\eta_u},
  \qquad
  \eta_o\ \ \text{in all remaining cases},
\]
the remaining cases being direction $C$ with $y\ge1$, direction $P$, and
direction $O$ with $y\ge1$.  The three are ordered by $k_1-K=(K-1)(\eta-1)\ge0$ and
$K\eta-k_1=\eta-1\ge0$, that is
\[
  \frac1{\eta_u}\ \le\ \frac{k_1}{K\eta_u}\ \le\ \frac{K\eta}{K\eta_u}=\eta_o .
\]
The outer two are attained, so $\max\tilde d/d=\eta_o$ and
$\max d/\tilde d=\eta_u$: the pair sits exactly on the band of
Definition~\ref{def:eta}, and its error is not merely bounded by $\eta$.

\emph{The optimum.}  $f(O)=f(0,0,K)=1$, and for every $S$ with $|S|\le K$ one
has $f(S)\le1$: at $y=K$ the value is exactly $1$, and at $y\le K-1$,
$f(S)=1-q^{x}\frac{K-y}{K}+z\delta_j\le1$ because
$z\delta_j\le(K-j)\frac{q^{j}}{K\eta}\le\frac{q^{x}}{K}\le q^{x}\frac{K-y}{K}$
whenever $y\le K-1$, using $z\le K-j\le\eta$ and $q^{j}\le q^{x}$.  So
$f(O^{\ast})=1$ with $O^{\ast}=O$.

\emph{The per-step gain table.}  The run is the following; the table lists, at
each state, the true and predicted gains of the three kinds of candidate.
\begin{center}
\begin{tabular}{llllll}
\toprule
step $t$ & state $(x,z,y)$ & candidate & true gain & predicted gain & picked \\
\midrule
$t<j$ & $(t,0,0)$ & $e\in C$ & $q^{t}/k_1$ & $q^{t}/(K\eta_u)$ & yes \\
      &           & $e\in P$ & $q^{\,j}/(K\eta)$ & $q^{\,j}/(K\eta_u)$ & no \\
      &           & $e\in O$ & $q^{t}/K$ & $q^{t}/(K\eta_u)$ & no (tie) \\
\midrule
$t\ge j$ & $(j,t-j,0)$ & $e\in P$ & $q^{\,j}/(K\eta)$ &
  $q^{\,j}/(K\eta_u)$ & yes \\
      &            & $e\in O$ & $q^{\,j}/K$ & $q^{\,j}/(K\eta_u)$ & no (tie)\\
\bottomrule
\end{tabular}
\end{center}
The entries follow from the gain formulas above: at $y=0$ a $C$-step has true
gain $q^{t}(1-q)=q^{t}/k_1$ and predicted gain
$q^{t}(1-q)W(0)=q^{t}W(0)/k_1=q^{t}/(K\eta_u)$; a $P$-step has predicted gain
$\eta_o\delta_j=q^{\,j}/(K\eta_u)$; an $O$-step at $y=0$ has predicted gain
$q^{x}/(K\eta_u)$ by \eqref{eq:W0-identity}.

\emph{The greedy induction.}  Claim: for $0\le t\le K$ the state after $t$
steps is $(x,z,y)=(\min(t,j),\ \max(t-j,0),\ 0)$.  It holds at $t=0$.  Assume
it at $t<j$.  By the table the maximum predicted gain at that state is
$q^{t}/(K\eta_u)$, attained by every remaining element of $C$ and every
element of $O$, while every element of $P$ scores $q^{\,j}/(K\eta_u)$, which is
strictly smaller because $q\in(0,1)$ and $j>t$.  Adversarial tie breaking picks
an element of $C$, so the state becomes $(t+1,0,0)$.  Now assume the claim at
$t$ with $j\le t<K$.  The state is $(j,t-j,0)$, $C$ is exhausted, and the
predicted gains of the remaining $P$ elements and of the $O$ elements are both
$q^{\,j}/(K\eta_u)$: the step is a tie and the adversary picks an element of
$P$, which is available because $t-j<K-j=|P|$.  After $K$ steps
$S^{K}=C\cup P$ and
\[
  f(S^{K})=1-q^{\,j}+(K-j)\,\delta_j
  =1-q^{\,j}+\frac{(K-j)q^{\,j}}{K\eta}
  =1-q^{\,j}\Bigl(1-\frac{K-j}{K\eta}\Bigr)=V_j(\eta),
\]
which is $V_j(\eta)f(O^{\ast})$ because $f(O^{\ast})=1$.
Choosing $j$ with $\eta\in[K-j,K-j+1]$ (that is $j$ the minimizing index of the
previous paragraph) gives $\rho_K(\eta)\le\min_iV_i(\eta)$, which together with
the lower bound is Theorem~\ref{thm:exact}.
% [HAND-PROOF-UNREVIEWED write-up; the identities are VERIFIED-SYMBOLIC for
% general K modulo one branch enumeration and VERIFIED-LP on the full lattice
% for K <= 6, results/N2_check.py 480/480.]

\emph{The role of the ties.}  Every one of the $K$ steps above is a tie with an
element of $O$, so adversarial tie breaking is load bearing.  On the steps
$t\ge j$ the tie is forced by the band: the prediction constraint of
\eqref{eq:redlp} holds with equality at the optimum, so
$\eta_o\,d_P=d_O/\eta_u$, and the largest value the band allows a $P$ element
and the smallest it allows an $O$ element coincide.  Under strict preference
the same value is an infimum rather than a maximum: for $\eta$ in the
interior of the segment $[K-j,K-j+1]$, replacing $\eta$ by
$\eta'<\eta$ inside $f$ while declaring the band at $\eta$ makes all
preferences strict and gives ratio $V_j(\eta')\to V_j(\eta)$ as
$\eta'\uparrow\eta$.
% The strict variant is checked at 6 (K, j, eta) triples and 2 values of
% epsilon in Part E of results/N2_check.py, all at segment INTERIOR points.
% K6 caveat (J2 review section 3.2): at a segment left endpoint the same
% perturbation can leave the family (exact rational counterexample at
% K=3, j=1, eta=2, eta'=19/10: one late O gain becomes -1/72, so f loses
% monotonicity; results/J2_extra_audit.py).  A repair sketch, approach the
% integer endpoint eta >= 2 from the left with the adjacent j+1 segment
% (adjacent branch values agree at breakpoints), and at eta = 1 perturb f
% and ftilde together by a vanishing positive modular term, is recorded
% here as [HAND-PROOF-UNREVIEWED] and not used by the theorem, whose
% statement is for adversarial ties.

The identities of this proof are machine-checked by
\texttt{results/N2\_check.py}.

% ---------------------------------------------------------------------------
\subsection{Structure of the attaining runs
(Proposition~\ref{prop:rigidity})}\label{app:rigidity}
% Status (ledger card T6b, H-J3).  The recurrences of Steps 3 to 6 and the
% rearrangement used in Step 7 are [VERIFIED-SYMBOLIC] and the collapse of the
% optimal face is [VERIFIED-LP], both in results/H_J3_gate_check.py; the
% assembly of the steps into a statement about a run of arbitrary size is
% [HAND-PROOF-UNREVIEWED] and is NOT to be upgraded (no script certifies an
% arbitrary K).  Source: results/J3_proof_structure.md sections 2, 2.1, 2.2.
% The J3 note's own oracle script (J3_structure_oracles.py) was never
% delivered, so every citation below names results/H_J3_gate_check.py, our
% independent re-implementation of that gate, and no claim rests on the
% undelivered script.  Per the file convention the visible text states the
% argument and never asserts a verification tag.

Throughout this subsection the notation is that of Appendix~\ref{app:exact},
$f(O^{\ast})=1$, and $\eta$ lies in the interior of a segment: either
$K-j<\eta<K-j+1$ for some $j$ with $1\le j\le K-1$, or $j=0$ and $\eta>K$.
Write $\gamma=1-\tfrac1\eta$ and $r_t=1-\sum_{s<t}d_s$, so that $r_0=1$ and
$r_{t+1}=r_t-d_t$.  Fix a run of predictive greedy with adversarial ties whose
value attains the exact worst case, $\sum_{t<K}d_t=V_j(\eta)$;
Appendix~\ref{app:validity} makes the vector $(d_t,g_{t,i})$ it induces
feasible for \eqref{eq:redlp}.  Nothing below assumes that the run avoids
$O^{\ast}$ or that the numbers $g_{t,i}$ are equal across $i$; both come out
of the argument.

\paragraph{Step 1: which multipliers are positive.}
On the interior of the segment, $M=K\eta-(K-j)>0$, $q\in(0,1)$, $k_1>0$, and
$\eta>1$: for $j\le K-2$ the segment gives $\eta>K-j\ge2$, and for $j=K-1$ the
segment is $(1,2)$.  Reading the signs off \eqref{eq:duals-jpos} and
\eqref{eq:duals-jzero}:
\begin{itemize}
\item $\lambda_S(t)>0$ exactly for $0\le t\le j$.  The entries with
$t\le j-1$ are positive multiples of $M$, and
$\lambda_S(j)=(K-j+1-\eta)/k_1>0$ because $\eta<K-j+1$ strictly inside the
segment, while $\lambda_S(t)=0$ for $t>j$ by definition.  For $j=0$ the single
entry is $\lambda_S(0)=1/\eta>0$.
\item $\lambda_C(t)>0$ exactly for $0\le t\le K-2$.  For $t\le j-1$ the entry
is a positive multiple of $M$; for $j\le t\le K-2$ it is
$(K-1-t)/(K(\eta-1))>0$; and $\lambda_C(K-1)=0$.
\item $\lambda_P(t)>0$ exactly for $j\le t\le K-1$.  It is $0$ for $t<j$ by
definition, and for $t\ge j$ its numerator is
$\eta-(K-t)\ge\eta-(K-j)>0$, again strictly inside the segment.
\item every multiplier of the family $\mathrm{mono}$ is $0$.
\end{itemize}
At an integer $\eta$ one of the two strict inequalities above becomes an
equality, which is the content of
Remark~\ref{rem:rigidity-breakpoints}: $\lambda_P(j)$ vanishes at the left
endpoint of the segment and $\lambda_S(j)$ at the right endpoint.

\paragraph{Step 2: every supported constraint holds with equality.}
Identity \eqref{eq:combination} reads $\sum_{t<K}d_t-V_j=\Sigma$ for every
vector, and $\Sigma$ is a sum of terms $\lambda_r\,\mathrm{slack}_r$ with
$\lambda_r\ge0$ throughout and $\mathrm{slack}_r\ge0$ at a feasible point.
The run attains the value, so the left-hand side is $0$, so every term
vanishes, so
\[
  \mathrm{sum}(t)\ (0\le t\le j),\qquad
  \mathrm{cons}(t,i)\ (0\le t\le K-2),\qquad
  \mathrm{pred}(t,i)\ (j\le t\le K-1)
\]
hold with equality, for every index $i$.  The constraints carrying a zero
multiplier, namely $\mathrm{sum}(t)$ for $t>j$, $\mathrm{cons}(K-1,i)$ and
every $\mathrm{mono}(t,i)$, stay valid but are not forced; Step 7 uses two of
them as inequalities.

\paragraph{Step 3: the geometric phase, $t<j$.}
Let $t<j$; the range is empty when $j=0$.  Both $\mathrm{sum}(t)$ and
$\mathrm{sum}(t+1)$ carry indices at most $j$, so
$\sum_ig_{t,i}=r_t$ and $\sum_ig_{t+1,i}=r_{t+1}$.  Summing the $K$ equalities
$\mathrm{cons}(t,i)$ over $i$, which is legitimate because
$t\le j-1\le K-2$,
\[
  \gamma\sum_ig_{t+1,i}-\sum_ig_{t,i}+Kd_t=0,
  \qquad\text{that is}\qquad
  \gamma\,r_{t+1}=r_t-Kd_t .
\]
Substituting $r_{t+1}=r_t-d_t$ gives $d_t(K-\gamma)=r_t(1-\gamma)=r_t/\eta$,
and $K-\gamma=K-1+\tfrac1\eta=k_1/\eta$, so
\[
  d_t=\frac{r_t}{k_1},\qquad r_{t+1}=r_t\Bigl(1-\frac1{k_1}\Bigr)=q\,r_t .
\]
With $r_0=1$ this gives $r_t=q^{\,t}$ and $d_t=q^{\,t}/k_1$ for every $t<j$,
and $r_j=q^{\,j}$.

\paragraph{Step 4: the switch at $t=j$.}
The index $j$ lies in the support of both $\lambda_S$ and $\lambda_P$, so
$\mathrm{pred}(j,i)$ holds with equality for every $i$, which is
$g_{j,i}=\eta\,d_j$, and $\mathrm{sum}(j)$ holds with equality, which is
$\sum_ig_{j,i}=r_j=q^{\,j}$ (for $j=0$ this is $r_0=1=q^{0}$).  Hence
$K\eta\,d_j=q^{\,j}$, that is
\[
  d_j=\frac{q^{\,j}}{K\eta},\qquad
  g_{j,i}=\frac{q^{\,j}}{K}\quad\text{for every }i .
\]
The $K$ numbers $g_{j,i}$ come out equal; this is a conclusion of the equality
case, not a symmetry imposed on the program, which treats the indices $i$
separately.

\paragraph{Step 5: the frozen phase, $j\le t\le K-2$.}
Both $\mathrm{pred}(t,i)$ and $\mathrm{pred}(t+1,i)$ hold with equality, so
$g_{t,i}=\eta d_t$ and $g_{t+1,i}=\eta d_{t+1}$, and $\mathrm{cons}(t,i)$
holds with equality, so
\[
  \gamma\,g_{t+1,i}=g_{t,i}-d_t=(\eta-1)d_t,
  \qquad\text{while}\qquad
  \gamma\,g_{t+1,i}=\gamma\eta\,d_{t+1}=(\eta-1)d_{t+1} .
\]
Since $\eta>1$ this forces $d_{t+1}=d_t$ and then $g_{t+1,i}=g_{t,i}$.  With
Step 4, $d_t=q^{\,j}/(K\eta)$ and $g_{t,i}=q^{\,j}/K$ for every $t$ with
$j\le t\le K-1$ and every $i$.

\paragraph{Step 6: the optimal marginals before the switch.}
For $t<j$ the equality $\mathrm{cons}(t,i)$ reads
$g_{t,i}=d_t+\gamma\,g_{t+1,i}$.  Going downwards from
$g_{j,i}=q^{\,j}/K$ and assuming $g_{t+1,i}=q^{\,t+1}/K$,
\[
  g_{t,i}=\frac{q^{\,t}}{k_1}+\gamma\,\frac{q^{\,t+1}}{K}
  =\frac{q^{\,t}}{K}\Bigl[\frac{K}{k_1}
    +\frac{\eta-1}{\eta}\cdot\frac{(K-1)\eta}{k_1}\Bigr]
  =\frac{q^{\,t}}{K}\cdot\frac{K+(K-1)(\eta-1)}{k_1}
  =\frac{q^{\,t}}{K},
\]
because $K+(K-1)(\eta-1)=(K-1)\eta+1=k_1$.  Downward induction gives
$g_{t,i}=q^{\,t}/K$ for every $t\le j$ and every $i$, one index $i$ at a time.

\paragraph{Step 7: the last state.}
The multiplier $\lambda_C(K-1)$ is $0$, so Step 2 does not force
$\mathrm{cons}(K-1,i)$; the constraint is valid all the same
(Lemma~\ref{lem:app-cons}), and so is $\mathrm{mono}(K-1,i)$
(Lemma~\ref{lem:app-mono}).  Write $d=d_{K-1}$, $g=g_{K-1,i}$ and
$h=g_{K,i}$.  Rearranged as the sharp form of Section~\ref{sec:coherence},
the two constraints read
\[
  d-\frac g\eta\;\ge\;\Bigl(1-\frac1\eta\Bigr)(g-h)
  \qquad\text{and}\qquad
  g-h\;\ge\;0 .
\]
Step 2 forces $\mathrm{pred}(K-1,i)$ with equality, that is $d=g/\eta$, so the
left-hand side of the first display is $0$; since $\eta>1$ the two displays
leave $h=g=q^{\,j}/K$.  Every variable of the run is therefore determined, and
the values are those of Proposition~\ref{prop:rigidity}.  As a check,
$\sum_{t<K}d_t=(1-q^{\,j})+(K-j)q^{\,j}/(K\eta)=V_j(\eta)$, and the gain table
of Appendix~\ref{app:exact} shows that the instances \eqref{eq:vj-family}
follow exactly this trajectory, so the forced trajectory is realized.

\paragraph{What the argument delivers and what it leaves open.}
Every $g_{K,i}$ equals $q^{\,j}/K>0$, while the zeroing convention of
Appendix~\ref{app:validity} sets $g_{K,i}=0$ whenever $o_i\in S^{K}$.  So no
optimal element is selected: an attaining run in this range of $\eta$ is
disjoint from the fixed $O^{\ast}$.  The reduction of
Appendix~\ref{app:validity} still has to admit runs that select optimal
elements, and the equality analysis excludes them only in the narrower
situation treated here.  The conclusion is about the vector $(d_t,g_{t,i})$
and about nothing else: the pair $(f,\tilde f)$ is not determined by it, and
neither are the true marginals of the candidates the run does not select.  In
particular the largest true marginal at a state, the quantity that
Definition~\ref{def:etasel} compares the chosen gain against, is not pinned by
this argument.  At an integer $\eta$ the argument stops at Step 1, and
Remark~\ref{rem:rigidity-breakpoints} records that the conclusion itself
fails there.

Two consequences of the same certificate are recorded in
Remark~\ref{rem:rigidity-ties}.  For the first, the slack of
$\mathrm{pred}(t,i)$ decomposes in case (b) of Appendix~\ref{app:validity},
with $P=\tilde d_{e_t}(S^{t})$ and $Q=\tilde d_{o_i}(S^{t})$, as
\[
  d_t-\frac{g_{t,i}}{\eta}
  =\underbrace{\frac{\eta_od_t-P}{\eta_o}}_{\text{upper band at }(S^{t},e_t)}
  +\underbrace{\frac{P-Q}{\eta_o}}_{\text{greedy}}
  +\underbrace{\frac{\eta_uQ-g_{t,i}}{\eta_u\eta_o}}_{\text{lower band at
   }(S^{t},o_i)} ,
\]
the three-term counterpart of the decomposition of the $\mathrm{cons}(t,i)$
slack in Lemma~\ref{lem:app-cons}.  Steps 2 and 7 put both slacks at zero,
the coherence one for $t\le K-2$ and the greedy-choice one at $t=K-1$, and
each term of each decomposition is nonnegative, so $P=Q$ at every step: the
selected element and every optimal element carry the same predicted gain, and
the run is a sequence of ties.  For the second, if the run outputs
$V_j(\eta)+\varepsilon$ instead of $V_j(\eta)$, then \eqref{eq:combination}
gives $\sum_r\lambda_r\,\mathrm{slack}_r=\varepsilon$ with all terms
nonnegative, hence $\mathrm{slack}_r\le\varepsilon/\lambda_r$ for every
constraint with $\lambda_r>0$.  That bound is stated for one run at a time;
the multipliers themselves degenerate at the segment endpoints, and no uniform
statement over $\eta$ and $K$ is derived from it here.

The recurrences of Steps 3 to 6, the rearrangement of Step 7 and the two
$K=3$, $\eta=2$ sequences of Remark~\ref{rem:rigidity-breakpoints} are
machine-checked by \texttt{results/H\_J3\_gate\_check.py}.  The same script
also fixes the objective of \eqref{eq:redlp} at the exact rational $V_j(\eta)$
and minimizes and maximizes each coordinate separately, for $K=2,\dots,6$,
every $j$ and two error levels per segment: on all $40$ optimal faces the
$2{,}480$ linear programs collapse every coordinate range to the values above,
with a largest deviation of about $4\cdot10^{-15}$, and no symmetry constraint
is added to the program.  What no script checks is the assembly of Steps 1 to
7 into a statement about a run of arbitrary size: the face enumeration is
finite, and the induction above carries the general case.

% ---------------------------------------------------------------------------
\subsection{The ceiling (Theorem~\ref{thm:ceiling})}\label{app:ceiling}
% Status: [HAND-PROOF-UNREVIEWED].  The night-4 draft flagged the randomized
% constant (1-K/n)/eta + K/n as needing a check; the averaging is written out
% below and the constant now comes out of the computation rather than being
% quoted.

Fix $\eta_u,\eta_o\ge1$, $K\ge1$, $n\ge2K$, and a constant $c>0$.

\paragraph{The construction.}
Let $\tilde f(S)=c|S|$, which does not depend on the hidden optimum, and for a
$K$-set $O$ with $B=N\setminus O$ let
\[
  f_O(S)=c\Bigl(\frac{|S\cap B|}{\eta_o}+\eta_u|S\cap O|\Bigr).
\]
$f_O$ is modular with positive coefficients, hence monotone and submodular
with $f_O(\emptyset)=0$.  Its single-element gains are $c/\eta_o$ on $B$ and
$c\eta_u$ on $O$, against the constant predicted gain $c$, so the ratios
$\tilde d/d$ take the two values $\eta_o$ (on $B$) and $1/\eta_u$ (on $O$), and
the error is exactly $(\eta_u,\eta_o)$ with both factors attained.  Because
both functions are modular, the gain of any pair $(A,B')$ is the sum of
single-element gains along any chain from $A$ to $A\cup B'$, so every
all-pairs ratio is a weighted average of the two single-element values and the
all-pairs error is exactly $(\eta_u,\eta_o)$ as well.  Since
$\eta_u\ge1\ge1/\eta_o$, the $K$ largest coefficients are those of $O$ and
$f_O(O^{\ast})=f_O(O)=c\eta_uK$.

\paragraph{Deterministic algorithms.}
The transcript of a deterministic algorithm querying $\tilde f$ is the same for
every $O$, so its output $T$, $|T|\le K$, is fixed before $O$ is chosen.  With
$n\ge2K$ there is a $K$-set $O\subseteq N\setminus T$, and for it
\[
  \frac{f_O(T)}{f_O(O^{\ast})}=\frac{c|T|/\eta_o}{c\eta_uK}
  \le\frac{cK/\eta_o}{c\eta_uK}=\frac{1}{\eta_u\eta_o}=\frac1\eta .
\]

\paragraph{Randomized algorithms.}
Let $O$ be uniform among $K$-subsets, independent of the random string.  The
output $T$ is again independent of $O$, and writing $m=|T|\le K$,
$\mathbb E|T\cap O|=mK/n$ and $\mathbb E|T\cap B|=m(1-K/n)$, so
\[
  \mathbb E\Bigl[\frac{f_O(T)}{f_O(O^{\ast})}\Bigr]
  =\frac{m(1-\tfrac Kn)/\eta_o+\eta_u mK/n}{\eta_uK}
  =\frac mK\Bigl[\frac{1-\tfrac Kn}{\eta_u\eta_o}+\frac Kn\Bigr]
  \le\Bigl(1-\frac Kn\Bigr)\frac1\eta+\frac Kn ,
\]
and $\min_O\le\mathbb E_O$ turns this into a statement about one instance.

\paragraph{The matching upper bound by exhaustive search.}
For any $S$, enumerate $S=\{s_1,\dots,s_m\}$ and telescope both functions
along the same chain $S_0=\emptyset\subset S_1\subset\dots\subset S_m=S$.  The
upper band applied to each of the $m$ terms gives
$\tilde f(S)=\sum_i\tilde d_{s_i}(S_{i-1})\le\eta_o\sum_id_{s_i}(S_{i-1})
=\eta_of(S)$, and the lower band gives $\tilde f(S)\ge f(S)/\eta_u$.  Let
$\hat S$ maximize $\tilde f$ over all sets of size at most $K$.  Then
\[
  f(\hat S)\ \ge\ \frac{\tilde f(\hat S)}{\eta_o}
  \ \ge\ \frac{\tilde f(O^{\ast})}{\eta_o}
  \ \ge\ \frac{f(O^{\ast})}{\eta_u\eta_o}=\frac{f(O^{\ast})}{\eta},
\]
the middle inequality by the choice of $\hat S$.  So the constant $1/\eta$ is
reached by an algorithm that is unrestricted in queries but makes no use of
structure, which is the sense in which the ceiling is met.
% GLOSSARY.md, "any algorithm": this whole subsection lives in range (i),
% unbounded queries.  The poly-query range (ii) is Theorem \ref{thm:hardness},
% whose constant is different; the two must not be quoted interchangeably.

No script checks the computations of this proof.

% ---------------------------------------------------------------------------
\subsection{Asymptotics (Corollary~\ref{cor:limit})}\label{app:asymptotics}

Fix $\eta\ge1$.  For $L_K$, the substitution $u=1/(\eta K)\in(0,1]$ and the
expansion $\ln(1-u)=-u-u^{2}/2-\dots$ give
$K\ln\bigl(1-\tfrac{1}{\eta K}\bigr)=-\tfrac1\eta+O(1/K)$, so
$L_K(\eta)\to1-e^{-1/\eta}$; moreover $(1-\tfrac{x}{K})^{K}$ increases in $K$
for $x\in(0,1]$, so the convergence of $L_K(\eta)$ is monotone from above.  For
$\rho_K=\min_jV_j$, fix $\eta$ and let $K$ grow.  The minimizing index is the
unique $j$ with $\eta\in[K-j,K-j+1]$, that is $j=K-\lfloor\eta\rfloor$ for
noninteger $\eta$ and either endpoint value at an integer $\eta$, so
$K-j$ stays bounded while $j\to\infty$.  With $k_1=(K-1)\eta+1$,
\[
  \ln q^{\,j}=j\ln\Bigl(1-\frac1{k_1}\Bigr)=-\frac{j}{k_1}+O\Bigl(\frac{j}{k_1^{2}}\Bigr),
  \qquad \frac{j}{k_1}=\frac{K-\lfloor\eta\rfloor}{(K-1)\eta+1}\longrightarrow\frac1\eta,
  \qquad \frac{j}{k_1^{2}}\longrightarrow0,
\]
so $q^{\,j}\to e^{-1/\eta}$, while the second factor of $V_j$ satisfies
$1-\frac{K-j}{K\eta}=1-\frac{\lfloor\eta\rfloor}{K\eta}\to1$.  Hence
$\rho_K(\eta)=V_j(\eta)\to1-e^{-1/\eta}$, the same limit as $L_K$, and by
Remark~\ref{rem:exact-gap} the two curves differ by $O(1/K)$ at fixed $\eta$.
% [HAND-PROOF-UNREVIEWED for the two limits above; the corresponding limit
% inside the hardness theorem, L_K(etahat) -> 1 - e^{-1/eta}, is
% [VERIFIED-SYMBOLIC] in results/N5_asymptotics.py checks (ii.a), (ii.b), with
% a 30-digit numeric cross-check at K = 10^6.]
% GAP LEFT OPEN AND MARKED: the clause "monotonically within each segment" of
% Corollary \ref{cor:limit} is derived above for L_K only (the standard
% monotonicity of (1-x/K)^K).  For rho_K the minimizing index j = K - floor(eta)
% moves with K, and no monotonicity argument in K is given here; the claim rests
% on the Richardson-extrapolated numerics of R8 and is [CONJECTURE] at that
% level of generality.  A human should either weaken the clause in
% sections/results.tex or supply the missing monotonicity step.

The limit identities of this proof are machine-checked, in the form used by
Theorem~\ref{thm:hardness}, by \texttt{results/N5\_asymptotics.py}.

% ---------------------------------------------------------------------------
\subsection{Bounded-query hardness (Theorem~\ref{thm:hardness})}
\label{app:hardness}
% Source after K5: results/J2_hardness_repair.tex (the drop-in supplement of
% the J2 external review), adopted after the independent recheck
% results/H3_j2_recheck.py.  It REPLACES the six-step write-up taken from
% results/N5_bounded_query_hardness.tex, whose delta closed form and Phi
% inverse are superseded.
% Status decomposition:
%  - the four-row edge table, the monotone row h_y, the extremal factors
%    sqrt(theta)A and sqrt(theta)B, the exact product Psi(theta) = theta A B,
%    the inflation delta = AB - 1, the first and second differences of the
%    true objective, the calibration H_{K,tau}(eta) = L_K(thetabar), the
%    comparison Psi <= Phi and both limits:
%    [VERIFIED-SYMBOLIC] (24 identities) and [VERIFIED-LP] (264 exact
%    rational count-grid instances, 57,728 edges, K = 2..12, every integer
%    tau = 1..K-1), results/J2_core_oracles.py.
%  - the counting bound, the canonical-transcript induction and the two
%    averagings: [HAND-PROOF-UNREVIEWED].  No script certifies an arbitrary
%    adaptive or randomized algorithm.  This tag is not to be upgraded.
%  - what the four rows retire: the delta closed form was [CONJECTURE] at
%    general (K, tau, eta) because the earlier derivation covered only the
%    balanced strip and left "no other pair of sets binds harder" to the LP.
%    The table below is exhaustive over the count grid, so that gap is
%    closed; what it also shows is that Phi is the minimal SYMMETRIC band and
%    not the pair's error, which is why the old "error exactly eta" reading
%    of the Phi calibration does not hold (J2 review section 6.2).
%  - the earlier hypothesis (H), Phi(1) <= eta, and its sufficient condition
%    K >= tau(1 + 2/ln eta) are not needed: Psi inverts in closed form.
%    The derivative typo in the old invertibility argument (a missing chain
%    factor K in the second branch of Phi) disappears with the branch.

In this subsection "information-theoretic" is used with the following meaning,
stated here at its first use: the bound depends only on the number and the
size of the queries that the algorithm makes and uses no computational
assumption, so nothing below is a statement about running time.

\paragraph{The family.}
Fix integers $1\le\tau<K<n$ and a real $\theta\ge1$, and write $a=a_\theta$,
$A=A_\theta$, $B=B_\theta$ for the constants of Section~\ref{sec:hardness}.
For a hidden $K$-set $O\subseteq N$ and $S\subseteq N$ put $x=|S\setminus O|$
and $y=|S\cap O|$, and define
\[
  F_O(S)=\theta^{-1/2}\Bigl[1-a^{x}\Bigl(1-\frac yK\Bigr)\Bigr],
  \qquad
  G_O(S)=\begin{cases}
    1-a^{\,x+y}, & y\le\tau,\\[2pt]
    1-a^{\,x+\tau}\dfrac{K-y}{K-\tau}, & y\ge\tau,
  \end{cases}
\]
together with the balanced-region profile $\hat G(s)=1-a^{s}$.  The two
branches of $G_O$ agree at $y=\tau$.  Both functions depend on $S$ only
through the pair $(x,y)$, so Lemma~\ref{lem:app-count} applies to them with
the partition $N=(N\setminus O)\cup O$.  The instance is $f=F_O$ and
$\tilde f=G_O$; the prefactor $\theta^{-1/2}$ places the symmetric split and
cancels from every ratio below.

\paragraph{Every edge of the count grid.}
On an edge with positive true gain put
$r=\Delta G_O/(\sqrt\theta\,\Delta F_O)$, the predicted-to-true gain ratio of
that edge.  Subtracting the two displays gives the following exhaustive table.
\begin{center}
\begin{tabular}{lll}
\toprule
Direction & Starting state & $r$ \\
\midrule
$x$ & $0\le y\le\tau$ & $a^{y}K/(K-y)$ \\
$x$ & $\tau\le y<K$   & $A$ \\
$y$ & $0\le y\le\tau-1$ & $a^{y}/\theta$ \\
$y$ & $\tau\le y<K$   & $A$ \\
\bottomrule
\end{tabular}
\end{center}
An $x$-edge starting at $y=K$ has zero true and zero predicted gain, and no
$y$-edge starts at $y=K$, so neither carries a band constraint.  The edge from
$y=\tau-1$ to $y=\tau$ is in the third row and the edge from $y=\tau$ to
$y=\tau+1$ is in the fourth, so the junction of the two branches needs no
separate extension argument, and the edges leaving the balanced strip are
already in rows two and four.  Every $x$ cancels, so the table is the same for
all $n>K$.

Write $h_y=a^{y}K/(K-y)$ for the first row.  Then
\[
  \frac{h_{y+1}}{h_y}-1=\frac{(\theta-1)K+y}{\theta K(K-y-1)}\;\ge\;0 ,
\]
so $1=h_0\le h_y\le h_\tau=A$ for $0\le y\le\tau$.  The third row
$a^{y}/\theta$ decreases from $1/\theta$ to $a^{\tau-1}/\theta=1/(\theta B)$.
Since $A\ge1$ and $B\ge1$, the largest and the smallest values of $r$ over the
whole grid are $A$ and $1/(\theta B)$, and both are attained on actual edges
because $n>K$ and $1\le\tau<K$.  In the notation of Definition~\ref{def:eta},
where the band reads $1/\eta_u\le\tilde d_e(S)/d_e(S)\le\eta_o$ and
$\tilde d_e(S)/d_e(S)=\sqrt\theta\,r$, the smallest admissible factors of the
unscaled pair are therefore
\[
  \eta_o^{\mathrm{act}}=\sqrt\theta\,A,\qquad
  \eta_u^{\mathrm{act}}=\sqrt\theta\,B,\qquad
  \eta^{\mathrm{act}}=\theta AB=\frac{\theta K-1}{K-\tau}=\Psi(\theta) ,
\]
and both factors are at least $1$, as Definition~\ref{def:eta} requires.

\paragraph{A prescribed split of the error.}
Replacing $\tilde f=G_O$ by $\beta G_O$ with $\beta=\eta_o/(\sqrt\theta A)$
multiplies every $r$ by $\beta$, so the largest predicted-to-true ratio becomes
$\eta_o$ and the smallest becomes $1/\eta_u$, with
$\eta_u\eta_o=\eta^{\mathrm{act}}$ unchanged.  The canonical profile $\hat G$
is multiplied by the same constant and remains a function of $|S|$ alone, so
the rescaling carries no information about the identity of $O$.  The choice
$\beta=\sqrt{B/A}$ returns the symmetric split
$\eta_u=\eta_o=\sqrt{\eta^{\mathrm{act}}}$.

\paragraph{The true objective and its optimum.}
Write $\bar F=\sqrt\theta\,F_O$, a function of $(x,y)$.  Its first differences
are
\[
  \Delta_x\bar F=a^{x}(1-a)\Bigl(1-\frac yK\Bigr),\qquad
  \Delta_y\bar F=\frac{a^{x}}{K},
\]
nonnegative on every feasible edge, and its second differences are
\[
  \Delta_x^{2}\bar F=-a^{x}(1-a)^{2}\Bigl(1-\frac yK\Bigr),\qquad
  \Delta_y\Delta_x\bar F=\Delta_x\Delta_y\bar F=-\frac{a^{x}(1-a)}{K},\qquad
  \Delta_y^{2}\bar F=0,
\]
all nonpositive.  Lemma~\ref{lem:app-count} then makes $F_O$ a normalized
monotone submodular function on every finite ground set.  Also
$0\le\bar F\le1$ and $\bar F(O)=1$, so $\max_{|S|\le K}F_O(S)=F_O(O)=\theta^{-1/2}$ and the optimal
set is $O^{\ast}=O$.  If $|T|\le K$ and $T\cap O=\emptyset$, then
\[
  \frac{F_O(T)}{F_O(O^{\ast})}=1-a^{|T|}\;\le\;1-a^{K}=L_K(\theta).
\]

\paragraph{Counting.}
For a fixed $S$ with $|S|\le K$ and a uniformly random $K$-subset $O$ of the
$n$-element ground set, the event $\{|S\cap O|>\tau\}$ is the union, over the
$(\tau+1)$-subsets $R\subseteq S$, of the events $\{R\subseteq O\}$, and
$\Pr[R\subseteq O]=\prod_{i=0}^{\tau}\frac{K-i}{n-i}\le(K/n)^{\tau+1}$, since
$(K-i)n\le K(n-i)$ for $n\ge K$.  A union bound over at most
$\binom{K}{\tau+1}\le K^{\tau+1}/(\tau+1)!$ sets $R$ gives
\[
  \Pr\bigl[|S\cap O|>\tau\bigr]
  \;\le\;\binom{K}{\tau+1}\Bigl(\frac Kn\Bigr)^{\tau+1}
  \;\le\;\frac{K^{2\tau+2}}{(\tau+1)!\,n^{\tau+1}} .
\]
The integrality of $\tau$ enters at this step, which is the reason for
$\tau=\lceil c\rceil+1$ rather than $\tau=c+1$ at real $c$.

\paragraph{The canonical transcript, and the deterministic statement.}
Let $\mathcal A$ be deterministic, make at most $Q=n^{c}$ queries of size at
most $K$, and return a set of size at most $K$.  Run it against the canonical
oracle $S\mapsto\hat G(|S|)$, rescaled by $\beta$ if a split was prescribed,
and let $S_1,\dots,S_Q$ and $T_0$ be the queries and the output it produces.
These are fixed sets, determined before $O$ is drawn.  Summing the counting
bound over the $Q$ canonical queries and adding
$\Pr[T_0\cap O\ne\emptyset]\le\mathbb E|T_0\cap O|=|T_0|K/n\le K^{2}/n$,
\[
  \Pr\bigl[\exists i:\ |S_i\cap O|>\tau\ \text{ or }\ T_0\cap O\ne\emptyset\bigr]
  \;\le\;\frac{K^{2\tau+2}}{(\tau+1)!\,n^{\tau+1-c}}+\frac{K^{2}}{n} .
\]
Under $n\ge4K^{c+2}$ the first term is at most
\[
  \frac{K^{\,c(c+1-\tau)}}{(\tau+1)!\,4^{\,\tau+1-c}}
  \;\le\;\frac1{16\,(\tau+1)!}\;\le\;\frac1{32},
\]
because $\tau\ge c+1$ makes the exponent of $K$ nonpositive and
$\tau+1-c\ge2$; the second term is $K^{2}/n\le1/(4K^{c})\le1/4$.  The sum is
at most $9/32$, so at least one $K$-set $O$ makes every canonical query
balanced and misses $T_0$; fix such an $O$.  Induction on $i$: if the first
$i-1$ answers of $G_O$ agree with the canonical ones, determinism makes the
$i$th query equal to $S_i$, and $|S_i\cap O|\le\tau$ puts $S_i$ in the first
branch, where $x+y=|S_i|$ and $G_O(S_i)=\hat G(|S_i|)$.  So all $Q$ answers
agree, the actual output is $T_0$, and $T_0\cap O=\emptyset$ gives
$f(T_0)/f(O^{\ast})\le1-a^{K}=L_K(\theta)$.

The order of the two arguments matters: the union bound is taken over the
canonical transcript, which does not depend on $O$.  A union bound over the
queries an algorithm actually issues against $G_O$ would not by itself make
those queries independent of $O$.

\paragraph{Randomized algorithms, by averaging.}
Fix the random seed $r$; then $\mathcal R_r$ is deterministic, with canonical
queries and canonical output $T_0^{(r)}$, and let $E_r$ be the event that all
of its canonical queries are balanced.  On $E_r$ the run against $G_O$ returns
$T_0^{(r)}$.  Every output has at most $K$ elements, so its counts satisfy
$x\le K$, and $a\in(0,1)$ gives
\[
  1-a^{x}\Bigl(1-\frac yK\Bigr)\;\le\;1-a^{K}+\frac yK ,
\]
while off $E_r$ the output ratio is at most $1$.  Hence, pointwise in $O$,
\[
  \frac{f_O(T)}{f_O(O^{\ast})}
  \;\le\;1-a^{K}+\frac{|T_0^{(r)}\cap O|}{K}+\mathbf 1[E_r^{c}] .
\]
Averaging over a uniform $O$ bounds the middle term by $K/n$ and the last one
by the query union bound above; averaging over $r$, exchanging the two finite
expectations and using $\min_O\le\mathbb E_O$ gives a single $O$ with the
additive $\varepsilon_n$ of Theorem~\ref{thm:hardness}.  Only the averaging
direction of the minimax principle is used; no minimax theorem is invoked.

\paragraph{Calibration, comparison with the symmetric band, and the limit.}
Choosing $\theta=\bar\theta=\Psi^{-1}(\eta)=(\eta(K-\tau)+1)/K$ makes the
pair's error exactly $\eta$ and turns the value bound into
\[
  L_K(\bar\theta)=1-\Bigl(1-\frac1{\bar\theta K}\Bigr)^{K}
  =1-\Bigl(1-\frac1{\eta(K-\tau)+1}\Bigr)^{K}=H_{K,\tau}(\eta),
\]
the constant of Theorem~\ref{thm:hardness}.  Since $\eta\tau\ge1$ gives
$\bar\theta\le\eta$ and $L_K$ is decreasing, $H_{K,\tau}(\eta)\ge L_K(\eta)$,
so Theorem~\ref{thm:hardness} and Proposition~\ref{prop:guarantee} do not
cross at any finite $K$.

For the same integer $\tau$, $\Psi(\theta)=\theta AB\le\theta\max\{A,B\}^{2}
=\Phi(\theta)$, with equality only at $A=B$.  Hence
$\bar\theta\ge\hat\eta=\Phi^{-1}(\eta)$ wherever the latter is defined, and
$H_{K,\tau}(\eta)=L_K(\bar\theta)\le L_K(\hat\eta)$: the constant above is not
weaker than the one the symmetric-band calibration produces.  The two differ
already at $(K,\tau,\theta)=(4,1,4)$, where the pair has factors
$(\eta_u,\eta_o)=(2,\tfrac52)$ and error $5$, while $\Phi(4)=\tfrac{25}4$; a
family calibrated through $\Phi^{-1}$ therefore has error at most, and in
general below, the prescribed level, which is what the earlier statement of
"error exactly $\eta$" did not account for.

The inflation relative to $\theta$ is $\delta(\theta)=AB-1=(\tau-1/\theta)/
(K-\tau)$, so $K\delta(\theta)\to\tau-1/\theta$ with $\tau$ and $\theta$
fixed.  With $c$ and $\eta$ fixed, $K/(\eta(K-\tau)+1)\to1/\eta$ and
$H_{K,\tau}(\eta)\to1-e^{-1/\eta}$.  No branch has to be inverted and no case
distinction enters the design parameter.

The edge table, the extremal factors, the exact product $\Psi$, the second
differences, the calibration identity and both limits are re-checked by
\texttt{results/J2\_core\_oracles.py}, symbolically for general parameters and
on an exact rational count grid; the extremes and the calibration comparison
are recomputed independently by \texttt{results/H3\_j2\_recheck.py}.  The
limit check of \texttt{results/N5\_asymptotics.py}, cited at the end of
Appendix~\ref{app:asymptotics}, is the corresponding statement for
$L_K(\hat\eta)$ and so belongs to the superseded calibration; the limit of
$H_{K,\tau}$ is the one checked in \texttt{results/J2\_core\_oracles.py}.  No
script certifies the counting bound, the transcript induction or the two
averagings: those are the hand-written part of this subsection, and finite
computations over instances do not range over adaptive algorithms.

% ---------------------------------------------------------------------------
\subsection{The greedy budget (Corollary~\ref{cor:greedybudget})}
\label{app:greedybudget}
% L1 (ledger card T10b).  The corollary is the family and the assembly of
% app:hardness at tau = 1, with the counting redone for the explicit budget
% Q = nK in place of n^c.  Status decomposition:
%  - every displayed inequality of the counting chain below, the calibration
%    thetabar >= 1, and H_{K,1} = U_K: [VERIFIED-SYMBOLIC],
%    results/L1_table.py section 1.
%  - the canonical-transcript induction and the two averagings are used
%    verbatim from app:hardness: [HAND-PROOF-UNREVIEWED], and per the ledger
%    this tag is NOT to be upgraded.

Fix $K\ge2$, $\eta>1$ and $n\ge4K^{5}$, and take the family of
Appendix~\ref{app:hardness} with $\tau=1$.  The constraint
$\eta\ge(K-1)/(K-\tau)$ of Theorem~\ref{thm:hardness} reads $\eta\ge1$ and
holds vacuously, and
$\bar\theta-1=(\eta-1)(K-1)/K\ge0$, so the calibration is admissible and the
value bound of the family is
$L_K(\bar\theta)=H_{K,1}(\eta)=U_K(\eta)$.  The edge table, the exact error
$\Psi(\theta)$, the prescribed-split rescaling and the optimum $O^{\ast}=O$
are exactly as in Appendix~\ref{app:hardness}; nothing in them refers to the
query budget.  Only the counting changes.

\paragraph{Counting at the budget $Q=nK$.}
At $\tau=1$ a query $S$, $|S|\le K$, is unbalanced when $|S\cap O|\ge2$,
which is the union over the pairs $R\subseteq S$ of the events
$\{R\subseteq O\}$.  For a uniformly random $K$-subset $O$,
\[
  \Pr[R\subseteq O]=\frac{K(K-1)}{n(n-1)}
  =\Bigl(\frac Kn\Bigr)^{2}-\frac{K(n-K)}{n^{2}(n-1)}
  \;\le\;\Bigl(\frac Kn\Bigr)^{2},
\]
so one query is unbalanced with probability at most
$\binom K2(K/n)^{2}$, and the canonical transcript's $Q=nK$ queries
$S_1,\dots,S_Q$ together with its output $T_0$ fail with probability at most
\[
  nK\binom K2\Bigl(\frac Kn\Bigr)^{2}+\Pr[T_0\cap O\ne\emptyset]
  \;\le\;\frac{K^{4}(K-1)}{2n}+\frac{K^{2}}{n}
  \;\le\;\frac{K^{5}}{2n}+\frac{K^{2}}{n}.
\]
Under $n\ge4K^{5}$ the first term is at most $\tfrac18$ and the second is at
most $1/(4K^{3})\le\tfrac1{32}$, so the total is at most
$\tfrac5{32}<\tfrac12$: some $K$-set $O$ makes every canonical query
balanced and misses $T_0$.  The transcript induction of
Appendix~\ref{app:hardness}, with balanced now meaning $|S_i\cap O|\le1$,
then forces the actual run to reproduce the canonical one and gives
$f(T_0)/f(O^{\ast})\le1-a^{K}=L_K(\bar\theta)=U_K(\eta)$, with error exactly
$\eta$ and any prescribed split realized by the $\beta$-rescaling.  The
ceiling $1/\eta$ holds for every deterministic algorithm at $n\ge2K$
(Theorem~\ref{thm:ceiling}), and $4K^{5}\ge2K$, so the two caps combine into
$\min\{U_K(\eta),1/\eta\}$ for the worst-case ratio.

The only property of the budget that the computation uses is
$Q\le n^{2}/(4K^{4})$, which $Q=nK$ satisfies exactly when $n\ge4K^{5}$.
By contrast, covering the budget $nK$ on the scale of
Theorem~\ref{thm:hardness} needs $n^{c}\ge nK$, hence $c>1$ and
$\tau=\lceil c\rceil+1\ge3$, and delivers only the weaker constant
$H_{K,3}(\eta)>H_{K,1}(\eta)$; running the count at the exact budget is what
keeps the $\tau=1$ constant.

\paragraph{Randomized algorithms.}
The averaging of Appendix~\ref{app:hardness}, applied verbatim with the
union bound above, yields for every randomized algorithm of the class a
single instance with
$\mathbb E[f(T)/f(O^{\ast})]\le U_K(\eta)+K/n+K^{5}/(2n)$; the corollary
states the slightly larger
$\varepsilon_n=K^{2}/n+K^{5}/(2n)$, which also matches the deterministic
output count.  The quantifier order is that of Theorem~\ref{thm:hardness}:
the instance may depend on the algorithm, and the expectation is over the
algorithm's randomness only.

\paragraph{The width of the interval.}
For the gap statement of Remark~\ref{rem:greedybudget}, expanding both
closed forms at fixed $\eta$ in $1/K$ (with $m=\lfloor\eta\rfloor$ and
$\rho_K=V_{K-m}$ on the active branch) gives
\[
  U_K(\eta)-\rho_K(\eta)
  =\frac{e^{-1/\eta}\,m\,(2\eta-m-1)}{2\eta^{2}}\cdot\frac1{K^{2}}
  +O\!\Bigl(\frac1{K^{3}}\Bigr):
\]
the constant and the $1/K$ coefficients of the two series agree (the latter
is the $c(\eta)$ of Appendix~\ref{app:asymptotics}), and the $1/K^{2}$
coefficients differ by $c'(\eta)$ as displayed.  At integer $\eta$ the
branches $m=\eta$ and $m=\eta-1$ of $c'$ agree, with common value
$e^{-1/\eta}(\eta-1)/(2\eta)$, so $c'$ is continuous although the
$1/K^{2}$ coefficient of $\rho_K$ alone is piecewise.  Both series, the
cancellation at order $1/K$, the closed form of $c'$ and its positivity for
$\eta>1$ are checked in \texttt{results/L1\_table.py} (section 2,
symbolically; section 3 numerically in exact rationals up to $K=400$,
with one Richardson step); the two ends of the interval are tabulated in
\texttt{results/L1\_table.csv}.  All of this is conditional on
Theorem~\ref{thm:exact} through the closed form of $\rho_K$.

% ---------------------------------------------------------------------------
\subsection{Queries of arbitrary size: a construction direction with
finite evidence}
\label{app:hardness-anysize}
% K9 (J4 audit): demoted from a theorem environment to a candidate bound.
% General-parameter admissibility of the family has only finite exhaustive
% support, so per the ledger card T12 this subsection presents a
% construction direction with finite evidence, not a theorem of the paper.
% Condensed from results/F3_hardness_full.tex (night-4 draft), which is the
% full write-up and is kept in the supplementary material.  Status tags are
% carried over from that file verbatim:
%   Step 1 (concentration, any query size)  [HAND-PROOF-UNREVIEWED]
%   Step 2 (the family is admissible)       [EXACT] 42 (K,eta) points,
%                                           results/N4_check.py section A;
%                                           [EXACT] 12 points,
%                                           results/F3_check.py section H
%   Step 3 (valuation / value)              [HAND-PROOF-UNREVIEWED], plus
%                                           [VERIFIED-SYMBOLIC] for the closed
%                                           form of W_K (N4_symbolic.py 19/19)
%   Step 4 (averaging, randomized)          [HAND-PROOF-UNREVIEWED]
%   The theorem as a whole                  [HAND-PROOF-UNREVIEWED]
%   j and m* as closed forms                [CONJECTURE], 264 points
%   D1, exponent 2 intrinsic                [VERIFIED-LP + EXACT], 48 configs
%   D2, two-sided band                      [VERIFIED-LP], 144 configs;
%                                           threshold n = K(T+1)+1 matched
%                                           at 12/12 points
%   Corollary (randomized) additionally uses F_O <= W_K + y/K, which no script
%   has checked.
% [EXACT] is that file's name for an exact-rational enumeration; it is used
% here unchanged so that the two documents can be compared line by line.
% NOTATION: the F3 draft writes T for the window constant j + m*, which
% collides with the output set T of Theorem \ref{thm:hardness}; it is written
% t* here.  Nothing else is renamed.

Theorem~\ref{thm:hardness} restricts every query to at most $K$ elements.  A
companion construction removes that restriction and pays for it in the budget.
For $K\ge3$ and $\eta>1$ put $k_1=\eta(K-1)+1$, $q=1-1/k_1$,
$\nu=\eta/(\eta-1)$, $j=\min\{K,\max\{0,K+1-\lceil\eta\rceil\}\}$,
$m^{\ast}=\lceil\eta K\rceil-1$ and $t^{\ast}=j+m^{\ast}$, and let
\[
  D=\max\Bigl\{\frac{q^{j}}{K\eta},\
    \frac{q^{j}\bigl(\nu^{m^{\ast}}/K-1\bigr)}
         {\eta(\nu^{m^{\ast}}-1)-m^{\ast}}\Bigr\},
  \qquad
  W_K(\eta)=1-q^{j}+(K-j)D .
\]
The family is the one of \texttt{results/F3\_hardness\_full.tex}: a hidden
$K$-set $O$, a true objective that is piecewise geometric in $x=|S\setminus O|$
and flattens at $x>t^{\ast}$, and a predictor that agrees with an $O$-free
profile $\hat G_{|S|}$ whenever $|S\cap O|\le1$ and also whenever
$|S\setminus O|>t^{\ast}$.  The second half of that condition is what removes
the query-size restriction: a query can break the simulation only if it meets
$O$ twice and has at most $t^{\ast}+K$ elements, so large queries need no
concentration argument at all.  Two caveats are recorded with the draft.  If
the tolerance is raised to $|S\cap O|\le\tau$ with $\tau\ge2$, the family is
infeasible at every inflation, by an explicit two-constraint certificate at
all $48$ tested configurations.  If instead the predictor is required to be a
function of $|S|$ on the whole two-sided band
$\bigl||S\cap O|-K|S|/n\bigr|\le\tau$, the family is feasible exactly when
$\tau=1$ and $n>K(t^{\ast}+1)$, and infeasible at every inflation otherwise;
which of the two formalizations a reader accepts is the open judgement call
of that draft.

\paragraph{Candidate bound (not a theorem of this paper).}
The intended statement is: for $K\ge3$, $\eta>1$ and $n\ge4K(t^{\ast}+K)$
(which forces $n\ge8K^{2}$), every deterministic algorithm
that makes at most $Q\le n^{2}/\bigl(2K^{2}(t^{\ast}+K)^{2}\bigr)$ queries to
$\tilde f$, of arbitrary size, and returns a set of size at most $K$, meets
a $K$-set $O$ for which the instance $(f,\tilde f)$ has monotone submodular
$f$, single-element error exactly $\eta$, $O^{\ast}=O$, and output
$T_{\mathcal A}$ with
\[
  \frac{f(T_{\mathcal A})}{f(O^{\ast})}\;\le\;W_K(\eta)
  \;=\;V_j(\eta)+(K-j)\Bigl(D-\frac{q^{j}}{K\eta}\Bigr),
\]
and, for randomized algorithms, the same bound in expectation up to an
additive $K/n+QK^{2}(t^{\ast}+K)^{2}/(2n^{2})$.  It is a candidate because
the admissibility of the family at general parameters is exactly what the
finite evidence below does not yet establish.

Since $t^{\ast}+K\le(2+\eta)K+1$, the admissible budget is
$Q=\Omega\bigl(n^{2}/((2+\eta)^{2}K^{4})\bigr)$: for fixed $K$ and $\eta$
this is a polynomial budget of degree approaching $2$, and not every
polynomial; a uniform reading such as $n^{2-o(1)}$ would additionally
require restricting how $K$ and $\eta$ grow with $n$, which is not done
here.  Within the $48$ checked configurations, the exponent $2$ is a
property of this family rather than of the analysis: it is the exponent in
$\Pr[|S\cap O|\ge2]\le(Ks/n)^{2}/2$ for a query of size $s$, and raising it
would require the $O$-oblivious window to tolerate $|S\cap O|\le\tau$ with
$\tau\ge2$, which is the first of the two regimes recorded above as
infeasible; no claim is made beyond those configurations or for extensions
of the family.
The constant is the reduced-LP segment value $V_j$ of
Theorem~\ref{thm:exact} plus an excess that decays quickly in $K$: at $\eta=2$
it is $3.2\cdot10^{-4}$, $5.9\cdot10^{-7}$, $4.5\cdot10^{-12}$ and
$5.2\cdot10^{-22}$ at $K=4,8,16,32$, and
$L_K(\eta)\le V_j(\eta)\le W_K(\eta)<U_K(\eta)$ with
$W_K(\eta)\to1-e^{-1/\eta}$.

What this supplementary statement does not carry is a general-$K$ argument.
The closed forms of $j$ and $m^{\ast}$ have numerical support at $264$
parameter points and no proof; monotonicity, submodularity and the
normalization of the family are enumerated exactly at $42$ $(K,\eta)$ points
and have no general-$K$ symbolic derivation; the concentration lemma, the
simulation induction, the assembly and the averaging are hand arguments that
no script checks, and the randomized version additionally uses an inequality
on the family that no script has checked.  The full draft, with its per-step
declarations, is \texttt{results/F3\_hardness\_full.tex} in the supplementary
material; the enumerations behind it are \texttt{results/N4\_check.py},
\texttt{results/F3\_check.py} and \texttt{results/N4\_symbolic.py}.

% ---------------------------------------------------------------------------
\subsection{Validity of the reduced constraints}\label{app:validity}
% Source: results/F2_R6_validity.tex, adapted to the paper notation (b_t -> e_t,
% O -> O*, the LP written as in \eqref{eq:redlp}).  STATUS after K4 (J2
% adoption): the slack identities behind pred(t,i) and cons(t,i) case (b) are
% [VERIFIED-SYMBOLIC] (results/H3_j2_recheck.py, results/J2_core_oracles.py),
% and the finite full-lattice slack minimization over TRUE lattice variables
% is [VERIFIED-LP] (J2_core_oracles.py: (n,K) = (4,2),(5,3), all optimal
% K-sets, four bands, 1,536 targets, min slack -6.7e-16).  The remaining
% bookkeeping (zeroing convention, case split, Lemma app-relaxation) is
% [HAND-PROOF-UNREVIEWED].
% Each family is written out one at a time, with the case e_t in O* (the gap
% flagged in R6) treated as a separate case in each.

This subsection supplies the step consumed by Appendix~\ref{app:exact}: every
constraint of \eqref{eq:redlp} is a valid inequality, so the linear program is
a relaxation of the set of realizable runs and its optimum is a lower bound on
the worst-case ratio for every $K$.

Let $f$ be monotone submodular with $f(\emptyset)=0$, let $\tilde f$ be an
arbitrary set function with $\tilde f(\emptyset)=0$ obeying the band of
Definition~\ref{def:eta}, and let $e_0,\dots,e_{K-1}$ and
$S^{0}=\emptyset,\dots,S^{K}$ be the picks and states of a run of predictive
greedy, so that $e_t\in\arg\max_{e\notin S^{t}}\tilde d_e(S^{t})$.  Fix an
optimal set $O^{\ast}=\{o_1,\dots,o_K\}$ with $f(O^{\ast})=1$ (if
$|O^{\ast}|<K$, set $g_{t,i}=0$ for the unused indices; every case (a) below
covers them).  Put $d_t=d_{e_t}(S^{t})$ and
\[
  g_{t,i}=\begin{cases} d_{o_i}(S^{t}), & o_i\notin S^{t},\\
                        0, & o_i\in S^{t}, \end{cases}
\]
the second line being the \emph{zeroing convention}, which is the device that
handles a step at which greedy selects an element of $O^{\ast}$: from that step
on the corresponding variable is $0$ rather than undefined.  It is consistent
with reading $g_{t,i}$ as a marginal gain, since an element already in $S^{t}$
contributes nothing further.  Finally $r_t=1-\sum_{s<t}d_s$, which telescopes
to $r_t=f(O^{\ast})-f(S^{t})$.

Throughout, $t\in\{0,\dots,K-1\}$, $i\in\{1,\dots,K\}$, and the three cases are
(a) $o_i\in S^{t}$; (b) $o_i\notin S^{t}$ and $o_i\ne e_t$; (c) $o_i=e_t$,
which is the case $e_t\in O^{\ast}$.

\begin{lemma}[signs]\label{lem:app-sign}
$d_t\ge0$ and $g_{t,i}\ge0$ for all $t,i$.
\end{lemma}
\begin{proof}
Each of these numbers is a marginal gain of $f$ or, under the zeroing
convention, equal to $0$.  This step uses monotonicity of $f$.
\end{proof}

\begin{lemma}[coverage, $\mathrm{sum}(t)$]\label{lem:app-sum}
$\sum_{i=1}^{K}g_{t,i}\ge r_t$ for every $t$.
\end{lemma}
\begin{proof}
Let $I_t=\{i:o_i\notin S^{t}\}$, so that
$\sum_ig_{t,i}=\sum_{i\in I_t}d_{o_i}(S^{t})$ by the zeroing convention; this
is exactly where a step with $e_s\in O^{\ast}$ for some $s<t$ is absorbed, the
indices of the already picked optimal elements dropping out of the sum.  Order
$I_t=\{i_1,\dots,i_m\}$ and telescope
\[
  f(S^{t}\cup O^{\ast})-f(S^{t})
  =\sum_{\ell=1}^{m}\Bigl[f\bigl(S^{t}\cup\{o_{i_1},\dots,o_{i_\ell}\}\bigr)
   -f\bigl(S^{t}\cup\{o_{i_1},\dots,o_{i_{\ell-1}}\}\bigr)\Bigr]
  \le\sum_{\ell=1}^{m}d_{o_{i_\ell}}(S^{t}),
\]
one application of submodularity per term.  Monotonicity gives
$f(S^{t}\cup O^{\ast})\ge f(O^{\ast})$, so
$\sum_ig_{t,i}\ge f(O^{\ast})-f(S^{t})=r_t$.
\end{proof}

\begin{lemma}[greedy choice inside the band, $\mathrm{pred}(t,i)$]
\label{lem:app-pred}
$d_t\ge g_{t,i}/\eta$ for every $t,i$.
\end{lemma}
\begin{proof}
Case (a): $g_{t,i}=0$ and $d_t\ge0$ by Lemma~\ref{lem:app-sign}.
Case (b): chain
\[
  d_t=d_{e_t}(S^{t})\ \ge\ \frac{\tilde d_{e_t}(S^{t})}{\eta_o}
  \ \ge\ \frac{\tilde d_{o_i}(S^{t})}{\eta_o}
  \ \ge\ \frac{d_{o_i}(S^{t})}{\eta_u\eta_o}=\frac{g_{t,i}}{\eta},
\]
the first inequality by the upper band at $(S^{t},e_t)$, the second by the
greedy choice ($o_i\notin S^{t}$ makes $o_i$ eligible at step $t$), the third
by the lower band at $(S^{t},o_i)$.
Case (c): here $g_{t,i}=d_{o_i}(S^{t})=d_{e_t}(S^{t})=d_t$ and the claim reads
$d_t\ge d_t/\eta$, which holds because $d_t\ge0$ and $\eta\ge1$.  No band
inequality is needed in this case, so the constraint does not degrade when
greedy happens to pick an optimal element.
\end{proof}

\begin{lemma}[single-element submodularity, $\mathrm{mono}(t,i)$]
\label{lem:app-mono}
$g_{t+1,i}\le g_{t,i}$ for every $t,i$.
\end{lemma}
\begin{proof}
Case (a): $o_i\in S^{t}\subseteq S^{t+1}$, so both sides are $0$.
Case (b): $o_i\notin S^{t+1}$ as well, so
$g_{t+1,i}=d_{o_i}(S^{t+1})\le d_{o_i}(S^{t})=g_{t,i}$ by submodularity in its
diminishing-returns form.
Case (c): $o_i=e_t\in S^{t+1}$, so $g_{t+1,i}=0$ while $g_{t,i}=d_t\ge0$ by
monotonicity.  This is the only place among the four families where the case
$e_t\in O^{\ast}$ replaces a submodularity argument by a monotonicity argument.
\end{proof}

\begin{lemma}[coherence, $\mathrm{cons}(t,i)$]\label{lem:app-cons}
$(1-\tfrac1\eta)g_{t+1,i}\ge g_{t,i}-d_t$ for every $t,i$.
\end{lemma}
\begin{proof}
Case (a): the left side is $0$ and the right side is $-d_t\le0$.
Case (b): abbreviate $d=d_t$, $g=g_{t,i}$, $h=g_{t+1,i}$ and, for the
predicted gains, $P=\tilde d_{e_t}(S^{t})$, $Q=\tilde d_{o_i}(S^{t})$,
$R=\tilde d_{o_i}(S^{t+1})$.  Because the four values of a set function on
$S^{t}$, $S^{t+1}=S^{t}\cup\{e_t\}$, $S^{t}\cup\{o_i\}$ and
$S^{t}\cup\{e_t,o_i\}$ determine both orders of insertion, the gains of
$e_t$ at $S^{t}\cup\{o_i\}$ are $d+h-g$ (true) and $P+R-Q$ (predicted).
The claim is then the sum of three nonnegative terms:
\[
  \Bigl(1-\frac1\eta\Bigr)h-g+d
  \;=\;\underbrace{\frac{\eta_o(d+h-g)-(P+R-Q)}{\eta_o}}_{\text{upper band
  at }(S^{t}\cup\{o_i\},\,e_t)}
  \;+\;\underbrace{\frac{P-Q}{\eta_o}}_{\text{greedy}}
  \;+\;\underbrace{\frac{\eta_uR-h}{\eta_u\eta_o}}_{\text{lower band at
  }(S^{t+1},\,o_i)} ,
\]
an identity in the six symbols; the first brace is nonnegative by
Definition~\ref{def:eta} at the off-trajectory state $S^{t}\cup\{o_i\}$, the
second because $o_i$ is eligible at step $t$ and $e_t$ was chosen, the third
by the lower band.  (Equivalently one may route the claim through part (ii)
of Lemma~\ref{lem:coherence}; the display above is the same argument with
its certificate written out.)
% The identity is [VERIFIED-SYMBOLIC]: results/H3_j2_recheck.py (sympy, this
% session's independent implementation) and results/J2_core_oracles.py (the
% external J2 review's implementation, plus 1,536 full-lattice LP targets
% minimizing this slack over true lattice variables, all nonnegative up to
% solver rounding).  The decomposition is due to the J2 external review
% (section 2), adopted in K4.
Case (c): $g_{t+1,i}=0$ by the zeroing convention and
$g_{t,i}-d_t=d_{e_t}(S^{t})-d_{e_t}(S^{t})=0$, so the constraint holds with
equality.  Lemma~\ref{lem:coherence}, which requires $e\ne e'$, is not invoked
in this case, which is what closes the gap this subsection exists to close.
\end{proof}

\begin{lemma}[relaxation]\label{lem:app-relaxation}
For every instance $(f,\tilde f)$ satisfying Definition~\ref{def:eta} with
$\eta=\eta_u\eta_o$, every optimal set $O^{\ast}$ and every run of predictive
greedy, the vector $(d_t,g_{t,i})$ is feasible for \eqref{eq:redlp} and its
objective value equals $f(S^{K})/f(O^{\ast})$.  Hence the optimum of
\eqref{eq:redlp} is a lower bound on the worst-case ratio of predictive greedy
at budget $K$ and error $\eta$.
\end{lemma}
\begin{proof}
Feasibility is Lemmas~\ref{lem:app-sign} to \ref{lem:app-cons}.  The objective
is $\sum_{t<K}d_t=f(S^{K})$ by telescoping, and $f(O^{\ast})=1$ by the
normalization.  Every realizable run gives a feasible point, so the minimum
over the feasible set is at most the minimum over realizable runs.
\end{proof}

\begin{remark}[what is and is not claimed]\label{rem:app-status}
Lemma~\ref{lem:app-relaxation} gives one direction only.  The reverse
direction, that every feasible point of \eqref{eq:redlp} is realized by some
instance, is not claimed here; the attaining instances of
Appendix~\ref{app:exact} supply what the theorem needs instead, and the
agreement of \eqref{eq:redlp} with the full-lattice program has been checked
for $K\le5$ and at nineteen $(K,\eta)$ test points.
\end{remark}

\begin{remark}[which error measure each family tolerates]\label{rem:app-rulers}
The families $\mathrm{sum}(t)$ and $\mathrm{mono}(t,i)$ use no error notion at
all.  The family $\mathrm{pred}(t,i)$ uses the band only at the greedy state
$S^{t}$, so it stays valid with $\eta$ replaced by $\etatr$, and, read
directly as a ratio of true gains at greedy states, with $\eta$ replaced by
$\etasel$.  The family $\mathrm{cons}(t,i)$ invokes
Lemma~\ref{lem:coherence}, whose proof uses the band at the state
$S^{t}\cup\{o_i\}$, which is off the greedy trajectory; so $\mathrm{cons}$ is
not covered by either restricted measure without a further assumption.  This
is the reason the $L_K$ bound, whose proof uses only $\mathrm{sum}$ and
$\mathrm{pred}$, is stated for $\etasel$ in
Proposition~\ref{prop:guarantee}, while the exact value of
Theorem~\ref{thm:exact} is stated for the global $\eta$.
\end{remark}

\begin{remark}[which hypothesis each family consumes]\label{rem:app-census}
$\mathrm{sum}(t)$ uses submodularity and monotonicity of $f$ and no property of
$\tilde f$.  $\mathrm{pred}(t,i)$ uses the greedy choice and both bands, and no
property of $f$ beyond monotonicity.  $\mathrm{mono}(t,i)$ uses submodularity,
and monotonicity in case (c).  $\mathrm{cons}(t,i)$ uses the greedy choice and
both bands through Lemma~\ref{lem:coherence}, and monotonicity in cases (a) and
(c).  No family uses submodularity of $\tilde f$, which is consistent with the
predictor being an arbitrary set function.
\end{remark}

The slack identities of $\mathrm{pred}$ and $\mathrm{cons}$ case (b) are
machine-checked by \texttt{results/H3\_j2\_recheck.py} and
\texttt{results/J2\_core\_oracles.py} (symbolic, plus 1{,}536 full-lattice
slack-minimization targets); the bookkeeping around them is not
oracle-checked.  The linear program the subsection justifies is compared
against the full-lattice program by \texttt{results/N3\_K5\_lattice.py}.

% ---------------------------------------------------------------------------
\subsection{Instances used in the experiments}\label{app:instances}
The worst-case instances run in Section~\ref{sec:experiments} are exactly the
families of Appendices~\ref{app:tightness} and \ref{app:exact}: the family
\eqref{eq:uk-family} at $\hat a=2$, and the family \eqref{eq:vj-family} for
$K\in\{3,5,8\}$ with one $\eta$ per segment $j$.  The experiment pipeline runs
them through the same lazy-greedy code as the real objectives, with the element
order implementing the adversarial tie direction ($B$ before $O$, and $C$ then
$P$ before $O$), and reproduces the values $L_K(\hat a)$ and $V_j(\eta)$ of the
two previous subsections to $10^{-10}$; on every $V_j$ instance it measures
$\etasel=\eta$, and on every member of \eqref{eq:uk-family} it measures
$\etasel=\etatr=\hat a<\eta$.

The computations of this subsection are machine-checked by
\texttt{results/E4\_worst\_instances.py} and
\texttt{results/F2\_etasel\_tight.py}.
